\setcounter{thema}{0}
\section{Θέματα Γ}
\def\xrwma{propAccent}
%=========== 1ο ΘΕΜΑ Γ ===========
\begin{thema}{Γ}
Δίνεται συνάρτηση $f:(0,+\infty)\to\mathbb{R}$ για την οποία ισχύει
\[ x^2f'(x)+xf(x)=1\ , \ x>0 \]
Να δείξετε ότι:
\begin{erwthma}
\item ο τύπος της $f$ είναι $f(x)=\dfrac{\ln{x}}{x}$ για κάθε $ x>0$.
\item υπάρχει τουλάχιστον ένα $x_0\in(0,1)$ τέτοιο ώστε $1-\ln{x_0}=2x_0^2$.

\item ισχύει $e\ln{x}\leq x$ για κάθε $x>0$.
\item Ένα σημείο $M$ κινείται πάνω στη γραφική παράσταση της $f$ του οποίου η τετμημένη αυξάνεται με ρυθμό $2\mu/s$. Τη στιγμή που το σημείο $M$ βρίσκεται στη θέση $A(e,f(e))$ να βρείτε το ρυθμό με τον οποίο μεταβάλλεται η τεταγμένη του.
\end{erwthma}
\end{thema}
%=========== ΤΕΛΟΣ - 1ο ΘΕΜΑ Γ ===========
%=========== 2ο ΘΕΜΑ Γ ===========
\begin{thema}{Γ}
Δίνεται παραγωγίσιμη συνάρτηση $f:\mathbb{R}\to\mathbb{R}$ με τύπο
\[ f(x)=\begin{cdcases}
x^2+a & ,x\geq0\\xe^x-x+\beta & ,x<0
\end{cdcases} \]
\begin{erwthma}
\item Να δείξετε ότι $a=\beta=-1$.
\item Να υπολογίσετε το όριο \[ \lim\limits_{x\to-\infty}{\frac{f(x)}{xf'(x)}} \]
\item Μελετήστε την $f$ ως προς την κυρτότητα και τα σημεία καμπής.
\item Να βρείτε την εξίσωση της εφαπτομένης της $C_f$ η οποία είναι κάθετη στην ευθεία $\zeta:y=x+1$.
\end{erwthma}
\end{thema}
%=========== ΤΕΛΟΣ - 2ο ΘΕΜΑ Γ ===========
%=========== 3ο ΘΕΜΑ Γ ===========
\begin{thema}{Γ}
Δίνεται η συνάρτηση $f:(0,+\infty)\to\mathbb{R}$ με τύπο $f(x)=x^{-ax}$ η οποία παρουσιάζει ακρότατο στη θέση $x=\frac{1}{e}$.
\begin{erwthma}
\item Να δείξετε ότι $a=1$.
\item Προσδιορίστε το είδος και την τιμή του ακρότατου και δείξτε ότι $ex\ln{x}+1\geq 0$ για κάθε $x\in(0,+\infty)$.
\item Να υπολογίσετε το όριο
\[ \lim\limits_{x\to+\infty}{\left(f(x)\cdot\hm{x}\right)} \]
\item Ένα σημείο $M$ ξεκινά από την αρχή των αξόνων και κινείται πάνω στη γραφική παράσταση της $f$ έτσι ώστε η τετμημένη του να αυξάνεται με ρυθμό $3\mu/\si{s}$. Να βρείτε το ρυθμό μεταβολής της τεταγμένης του $M$ τη χρονική στιγμή που αυτό διέρχεται από το $A\left(e,f(e)\right)$.
\end{erwthma}
\end{thema}
%=========== ΤΕΛΟΣ - 3ο ΘΕΜΑ Γ ===========

%=========== 4ο ΘΕΜΑ Γ ===========
\begin{thema}{Γ}
Δίνεται η παραγωγίσιμη συνάρτηση $f: \mathbb{R} \to \mathbb{R}$ για την οποία ισχύει $f(x)f'(x) = e^{2x}$ για κάθε $x \in \mathbb{R}$ και $f(0) = \sqrt{2}$. 
\begin{erwthma}
  \item Να αποδείξετε ότι $f^2(x) = e^{2x} + 1$ και στη συνέχεια να βρείτε τον τύπο της $f(x)$.
  \item Να αποδείξετε ότι η συνάρτηση είναι γνησίως αύξουσα στο $\mathbb{R}$ και να βρείτε το σύνολο τιμών της.
  \item Να ορίσετε την αντίστροφη συνάρτηση $f^{-1}$.
  \item Έστω ένα σημείο $M(x, y)$ που κινείται πάνω στη γραφική παράσταση της $f$. Αν η τετμημένη του αυξάνεται με σταθερό ρυθμό $2 \mu/\si{sec}$, να βρείτε τον ρυθμό μεταβολής του εμβαδού του τριγώνου $OAM$ τη χρονική στιγμή που το $M$ διέρχεται από τον άξονα $y'y$, όπου $O(0,0)$ και $A(x, 0)$ η προβολή του $M$ στον $x'x$.
\end{erwthma}
\end{thema}
%=========== ΤΕΛΟΣ - 4ο ΘΕΜΑ Γ ===========

%=========== 5ο ΘΕΜΑ Γ ===========
\begin{thema}{Γ}
Μια συνάρτηση $f: [0, +\infty) \to \mathbb{R}$ είναι παραγωγίσιμη με $f'(x) = \dfrac{1}{x^2+1}$ και $f(0) = 0$.
\begin{erwthma}
  \item Να μελετήσετε την $f$ ως προς την κυρτότητα και να βρείτε την εξίσωση της εφαπτομένης της $C_f$ στο σημείο $A(0, f(0))$.
  \item Να αποδείξετε ότι $f(x) \le x$ για κάθε $x \ge 0$ και να υπολογίσετε το όριο $\lim\limits_{x \to 0^+} \dfrac{f(x) - x}{x^2}$.
  \item Να αποδείξετε ότι $f(x+1) - f(x) < \dfrac{1}{x^2+1}$ για κάθε $x > 0$.
  \item Να αποδείξετε ότι η εξίσωση $f(x) = x - 1$ έχει ακριβώς μία ρίζα στο $(0, +\infty)$.
\end{erwthma}
\end{thema}
%=========== ΤΕΛΟΣ - 5ο ΘΕΜΑ Γ ===========

%=========== 6ο ΘΕΜΑ Γ ===========
\begin{thema}{Γ}
Δίνεται η συνάρτηση $f(x) = \ln x$ με πεδίο ορισμού $(0, +\infty)$.
\begin{erwthma}
\item Να αποδείξετε ότι υπάρχει $\xi \in (x, x+1)$ τέτοιο ώστε $\ln(x+1) - \ln x = \dfrac{1}{\xi}$.
\item Να αποδείξετε την ανισότητα $\dfrac{1}{x+1} < \ln\left(1 + \dfrac{1}{x}\right) < \dfrac{1}{x}$ για κάθε $x > 0$.
\item Με βάση την παραπάνω ανισότητα, αποδείξτε ότι $\lim\limits_{x \to +\infty} \left[x \ln\left(1 + \dfrac{1}{x}\right)\right] = 1$.
\item Να αποδείξετε ότι \[\ln{\frac{e+1}{2}}<\int_{1}^{e}{f\left(1+\frac{1}{x}\right)\d x}<1\]
\end{erwthma}
\end{thema}
%=========== ΤΕΛΟΣ - 6ο ΘΕΜΑ Γ ===========

%=========== 7ο ΘΕΜΑ Γ ===========
\begin{thema}{Γ}
Δίνεται οδικός χάρτης όπου η γραφική παράσταση της συνάρτησης $f(x) = \sqrt{x}, x \ge 0$ παριστάνει έναν δρόμο. Ένα όχημα $M$ κινείται πάνω σε αυτή την καμπύλη με τέτοιο τρόπο ώστε η τετμημένη του $x(t)$ να αυξάνεται με ρυθμό $1 \mu/\si{sec}$. To όχημα ρίχνει φως προς τον άξονα $x'x$ με το οποίο σχηματίζει μια εφαπτόμενη ευθεία.
\begin{erwthma}
  \item Να γράψετε τη γενική εξίσωση της εφαπτομένης της $C_f$ στο σημείο $M(x_0, f(x_0))$ και να βρείτε πού τέμνει τον άξονα $x'x$ (έστω σημείο $B$).
  \item Να δείξετε ότι το εμβαδόν του τριγώνου $O B M$, όπου $O(0,0)$ είναι η αρχή των αξόνων, είναι $E(x_0) = \dfrac{x_0\sqrt{x_0}}{2}$.
  \item Τη χρονική στιγμή $t_0$ που το όχημα περνάει από το σημείο με τετμημένη $4$, να βρείτε το ρυθμό μεταβολής της γωνίας $\theta$ που σχηματίζει η εφαπτομένη στο $M$ με τον άξονα $x'x$.
  \item Να υπολογίσετε τον αντίστοιχο ρυθμό μεταβολής του εμβαδού του τριγώνου $O B M$ την ίδια χρονική στιγμή $t_0$.
\end{erwthma}
\end{thema}
%=========== ΤΕΛΟΣ - 7ο ΘΕΜΑ Γ ===========

%=========== 8ο ΘΕΜΑ Γ ===========
\begin{thema}{Γ}
Έστω ότι για μια συνάρτηση $f: \mathbb{R} \to \mathbb{R}$, παραγωγίσιμη στο $\mathbb{R}$, ισχύει η ανισότητα $e^x - 1 \le f(x) \le x e^x + x$ για κάθε $x \ge 0$.
\begin{erwthma}
  \item Να αποδείξετε ότι $f(0) = 0$ και να υπολογίσετε την τιμή $f'(0)$.
  \item Αν $f(x) = e^x - 1$, να βρείτε, αν υπάρχει, την οριζόντια ασύμπτωτη της γραφικής παράστασης της συνάρτησης.
  \item Να αποδείξετε ότι ισχύει $xe^x + x \ge e^x - 1$ για κάθε $x \ge 0$.
  \item Να βρείτε το εμβαδόν του χωρίου που περικλείεται από τις γραφικές παραστάσεις των συναρτήσεων $g(x) = xe^x + x, h(x) = e^x - 1$ και των κατακόρυφων ευθειών $x=0$ και $x=1$.
\end{erwthma}
\end{thema}
%=========== ΤΕΛΟΣ - 8ο ΘΕΜΑ Γ ===========

%=========== 9ο ΘΕΜΑ Γ ===========
\begin{thema}{Γ}
Δίνεται η συνάρτηση $f(x) = e^x - \dfrac{x^2}{2} - x - 1$.
\begin{erwthma}
  \item Να μελετήσετε την $f$ ως προς την κυρτότητα.
  \item Να προσδιορίσετε το πρόσημό της $f'(x)$ σε όλο το $\mathbb{R}$.
  \item Να βρείτε το ελάχιστο της συνάρτησης $f$ και να αποδείξετε την ανισότητα $e^x \ge \dfrac{x^2}{2} + x + 1$ για κάθε $x \in \mathbb{R}$.
  \item Να αποδείξετε ότι $\dint_0^1 e^x \d x > \dfrac{5}{3}$.
\end{erwthma}
\end{thema}
%=========== ΤΕΛΟΣ - 9ο ΘΕΜΑ Γ ===========

%=========== 10ο ΘΕΜΑ Γ ===========
\begin{thema}{Γ}
Δίνεται η συνάρτηση $f(x) = \dfrac{x^2 - a x + 2}{e^x}$ όπου $a \in \mathbb{R}$. Γνωρίζουμε ότι η γραφική παράσταση της $f$ παρουσιάζει σημείο καμπής στο σημείο με τετμημένη $x_0 = 1$.
\begin{erwthma}
  \item Να αποδείξετε ότι $a = 1$.
  \item Να μελετήσετε την $f$ ως προς τη μονοτονία και να βρείτε τις θέσεις των τοπικών ακροτάτων.
  \item Να συγκρίνετε τους αριθμούς $f(2)$ και $f(3)$. 
  \item Να βρείτε τα όρια της $f$ στο $+\infty$ και $-\infty$ και να εξετάσετε την ύπαρξη ασύμπτωτων.
\end{erwthma}
\end{thema}
%=========== ΤΕΛΟΣ - 10ο ΘΕΜΑ Γ ===========

%=========== 11ο ΘΕΜΑ Γ ===========
\begin{thema}{Γ}
Δίνεται η συνάρτηση $f(x) = x^3 - 3x + 1$ η οποία παρουσιάζει δύο τοπικά ακρότατα στα σημεία $x_1$ και $x_2$.
\begin{erwthma}
  \item Να βρείτε τα σημεία αυτά, καθώς και τη μονοτονία της συνάρτησης.
  \item Να αποδείξετε ότι η εξίσωση $f(x) = 0$ έχει ακριβώς τρεις πραγματικές ρίζες $\rho_1 \in (-\infty, x_1), \rho_2 \in (x_1, x_2)$ και $\rho_3 \in (x_2, +\infty)$.
  \item Να προσδιορίσετε το σημείο καμπής της $C_f$ και να βρείτε την εξίσωση της εφαπτομένης της σε αυτό το σημείο.
  \item Να υπολογίσετε το εμβαδόν του χωρίου που περικλείεται από τη γραφική παράσταση της $f$, την εφαπτομένη της στο σημείο καμπής και την κατακόρυφη ευθεία $x = 1$.
\end{erwthma}
\end{thema}
%=========== ΤΕΛΟΣ - 11ο ΘΕΜΑ Γ ===========

%=========== 12ο ΘΕΜΑ Γ ===========
\begin{thema}{Γ}
Έστω η παραγωγίσιμη συνάρτηση $f: [1, e] \to \mathbb{R}$ για την οποία ισχύει $f(1)=0$ και $f(e)=e$.
\begin{erwthma}
  \item Να αποδείξετε ότι υπάρχει τουλάχιστον ένα σημείο $\xi \in (1, e)$ τέτοιο ώστε $f'(\xi) = \dfrac{e}{e-1}$.
  \item Να αποδείξετε ότι $f'(\xi) > 1$. Στη συνέχεια, αν $\min f'(x) = f'(\xi)$, τι συμπεραίνετε για την κλίση της εφαπτομένης στο $(1, e)$;
  \item Αν επιπλέον γνωρίζετε ότι $f(x) \ge \ln x$ για κάθε $x \in [1, e]$, να δείξετε ότι $f'(1) \ge 1$.
  \item Να αποδείξετε ότι υπάρχει τουλάχιστον ένα $x_0 \in (1, e)$ τέτοιο ώστε $f(x_0) = x_0 - 1$. 
\end{erwthma}
\end{thema}
%=========== ΤΕΛΟΣ - 12ο ΘΕΜΑ Γ ===========

%=========== 13ο ΘΕΜΑ Γ ===========
\begin{thema}{Γ}
Μια επιχείρηση κατασκευάζει κουτιά σχήματος ορθογωνίου παραλληλεπιπέδου χωρίς καπάκι, με τετράγωνη βάση πλευράς $x$ και ύψος $y$. Το συνολικό εμβαδόν επιφάνειας που χρησιμοποιείται είναι σταθερό και ισοδυναμεί με $12\ \si{m^2}.$
\begin{erwthma}
  \item Να αποδείξετε ότι ο όγκος $V$ του κουτιού συναρτήσει της πλευράς της βάσης $x$ δίνεται από τον τύπο $V(x) = 3x - \dfrac{x^3}{4}$ για $x \in (0, \sqrt{12})$.
  \item Να βρείτε τις διαστάσεις του κουτιού οι οποίες μεγιστοποιούν τον όγκο του. 
  \item Στη βέλτιστη κατασκευή, τι σχέση συνδέει την πλευρά της βάσης με το ύψος;
  \item Αν το κουτί που κατασκευάζεται έχει το μέγιστο όγκο, υπολογίστε τα όρια $\lim\limits_{x \to 2} \dfrac{V(x) - V(2)}{x-2}$ και $\lim\limits_{x \to \sqrt{12}^-} \dfrac{V(x)}{\sqrt{12}-x}$.
\end{erwthma}
\end{thema}
%=========== ΤΕΛΟΣ - 13ο ΘΕΜΑ Γ ===========

%=========== 14ο ΘΕΜΑ Γ ===========
\begin{thema}{Γ}
Δίνεται η συνάρτηση $f(x) = x \ln x - x + 1$ με $x > 0$.
\begin{erwthma}
  \item Να μελετήσετε την $f$ ως προς τη μονοτονία και τα ακρότατα.
  \item Να βρείτε την εφαπτομένη της γραφικής παράστασης της $f$ στο σημείο της με τετμημένη $1$.
  \item Αφού μελετήσετε την κυρτότητα της $f$, να αποδείξετε ότι ισχύει $x \ln x \ge x - 1$ για κάθε $x > 0$.
  \item  Να αποδείξετε ότι $\ln x < (x+1)\ln(x+1) - x\ln x - 1 < \ln(x+1)$ για κάθε $x > 0$.
\end{erwthma}
\end{thema}
%=========== ΤΕΛΟΣ - 14ο ΘΕΜΑ Γ ===========

%=========== 15ο ΘΕΜΑ Γ ===========
\begin{thema}{Γ}
Δίνεται μια συνάρτηση $f: \mathbb{R} \to \mathbb{R}$ η οποία ικανοποιεί τη σχέση $f(x+y) = f(x) + f(y)$ για κάθε $x, y \in \mathbb{R}$. Υποθέτουμε ότι η $f$ είναι παραγωγίσιμη στο $x_0=0$ με $f'(0) = 2$.
\begin{erwthma}
  \item Να αποδείξετε ότι $f(0) = 0$ και να γράψετε τον ορισμό της παραγώγου $f'(0)$ μέσω ορίου.
  \item Να αποδείξετε ότι η $f$ είναι παραγωγίσιμη στο $\mathbb{R}$ με $f'(x) = 2$ για κάθε $x \in \mathbb{R}$, και στη συνέχεια να βρείτε τον τύπο της $f$.
  \item Να υπολογίσετε το εμβαδόν του χωρίου που περικλείεται από τη γραφική παράσταση της $f$, της αντίστροφής της $f^{-1}$, και των κατακόρυφων ευθειών $x=2$ και $x=4$.
\end{erwthma}
\end{thema}
%=========== ΤΕΛΟΣ - 15ο ΘΕΜΑ Γ ===========

%=========== 16ο ΘΕΜΑ Γ ===========
\begin{thema}{Γ}
Για μια παραγωγίσιμη συνάρτηση $f(x)$ που ορίζεται στο $(0, +\infty)$ ισχύει η σχέση $x f'(x) - f(x) = x^2 e^x$ για κάθε $x > 0$ με $f(1) = e$.
\begin{erwthma}
  \item Να αποδείξετε ότι $f(x) = x e^x$.
  \item Να βρείτε τις ρίζες και το πρόσημο της $f(x)$, καθώς και τις ασύμπτωτες στο $+\infty$ και στο $0$.
  \item Να υπολογίσετε το όριο $\lim\limits_{x \to 0^+} \dfrac{f(x) - x}{x^2}$.
  \item Να βρείτε την εξίσωση της εφαπτομένης της $C_f$ η οποία διέρχεται από την αρχή των αξόνων $O(0,0)$.
\end{erwthma}
\end{thema}
%=========== ΤΕΛΟΣ - 16ο ΘΕΜΑ Γ ===========

%=========== 17ο ΘΕΜΑ Γ ===========
\begin{thema}{Γ}
Έστω η παραγωγίσιμη συνάρτηση $f: [0, 1] \to \mathbb{R}$ με $f(0) = 0$ και $f(1) = 1$. Υποθέτουμε ότι η πρώτη της παράγωγος είναι γνησίως αύξουσα στο $(0,1)$.
\begin{erwthma}
  \item Να αποδείξετε ότι υπάρχει $\xi \in (0, 1)$ τέτοιο ώστε $f'(\xi) = 1$.
  \item Να αποδείξετε ότι η τιμή $\xi$ είναι μοναδική και στη συνέχεια να βρείτε το πρόσημο της συνάρτησης $h'(x) = f'(x) - 1$ εκατέρωθεν του $\xi$.
  \item Να αποδείξετε την ανισότητα $f(x) < x$ για κάθε $x \in (0, 1)$.
  \item Να δείξετε ότι $\dint_0^1 f(x)\d x < \dfrac{1}{2}$.
\end{erwthma}
\end{thema}
%=========== ΤΕΛΟΣ - 17ο ΘΕΜΑ Γ ===========

%=========== 18ο ΘΕΜΑ Γ ===========
\begin{thema}{Γ}
Έστω η συνάρτηση $f(x) = e^x + x - 2$.
\begin{erwthma}
  \item Δείξτε ότι η εξίσωση $f(x) = 0$ έχει ακριβώς μία ρίζα στο διάστημα $x_0\in(0, 1)$.
  \item Να βρείτε, εφόσον υπάρχει, την αντίστροφη συνάρτηση $f^{-1}$.
  \item Να υπολογίσετε το όριο $\lim\limits_{x \to 0} \dfrac{f(x) + 1}{x}$.
  \item Αν η γραφική παράσταση της $f^{-1}$ διέρχεται από το σημείο $(-1, 0)$, να ορίσετε την εφαπτομένη της $C_f$ στο $(0, -1)$ και να εξηγήσετε πώς προκύπτει συμμετρικά η εφαπτομένη της αντίστροφης.
\end{erwthma}
\end{thema}
%=========== ΤΕΛΟΣ - 18ο ΘΕΜΑ Γ ===========

%=========== 19ο ΘΕΜΑ Γ ===========
\begin{thema}{Γ}
Δίνεται η συνάρτηση $f(x) = \dfrac{\ln x}{x}$ για $x > 0$.
\begin{erwthma}
  \item Να μελετήσετε την $f$ ως προς τη μονοτονία και να βρείτε το ολικό της ακρότατο.
  \item Συγκρίνετε τους αριθμούς $e^\pi$ και $\pi^e$.
  \item Να μελετήσετε την $f$ ως προς την κυρτότητα και τα σημεία καμπής της.
  \item Έστω $F(x)$ μια αρχική της $f(x)$ στο διάστημα $(0, +\infty)$. Να υπολογίσετε το $\lim\limits_{x \to +\infty} \dfrac{F(x)}{x^2}$.
\end{erwthma}
\end{thema}
%=========== ΤΕΛΟΣ - 19ο ΘΕΜΑ Γ ===========

%=========== 20ο ΘΕΜΑ Γ ===========
\begin{thema}{Γ}
Δίνεται η συνάρτηση $f(x) = \sqrt{x+1}$ με $x \ge -1$. 
\begin{erwthma}
  \item Να ορίσετε την εξίσωση της εφαπτομένης της $C_f$ στο σημείο με τετμημένη $x_0=0$.
  \item Να μελετήσετε τη συνάρτηση ως προς την κυρτότητα. Αποδείξτε στη συνέχεια, ότι $\sqrt{x+1} \le 1 + \dfrac{x}{2}$ για κάθε $x \ge -1$.
  \item Να αποδείξετε ότι $\dint_0^3 \sqrt{x+1} \d x \le \dfrac{15}{4}$.
  \item Να υπολογίσετε το όριο
\[\lim\limits_{x\to 0}{\dfrac{1}{2\sqrt{x+1}-2-x}}\]
\end{erwthma}
\end{thema}
%=========== ΤΕΛΟΣ - 20ο ΘΕΜΑ Γ ===========

%=========== 21ο ΘΕΜΑ Γ ===========
\begin{thema}{Γ}
Μια συνάρτηση $f: \mathbb{R} \to \mathbb{R}$ ικανοποιεί τη σχέση $f^3(x) + 3f(x) = x$ για κάθε $x \in \mathbb{R}$.
\begin{erwthma}
  \item Να αποδείξετε ότι η συνάρτηση $f$ είναι γνησίως αύξουσα. 
  \item Να επιλύσετε την εξίσωση $f(x) = 0$ και να δείξετε ότι η συνάρτηση $f(x)$ διατηρεί το ίδιο πρόσημο με το $x$.
  \item Να αποδείξετε ότι η συνάρτηση αντιστρέφεται και να υπολογίσετε τον τύπο της $f^{-1}(x)$.
  \item Υπολογίστε το ολοκλήρωμα $\dint_0^4 f(x)\d x$.
\end{erwthma}
\end{thema}
%=========== ΤΕΛΟΣ - 21ο ΘΕΜΑ Γ ===========

%=========== 22ο ΘΕΜΑ Γ ===========
\begin{thema}{Γ}
Δίνεται η συνάρτηση $f(x) = x \ln\left(\dfrac{x+1}{x}\right)$.
\begin{erwthma}
  \item Να βρείτε το πεδίο ορισμού και τις ασύμπτωτες της γραφικής παράστασης της $f$.
  \item Να μελετήσετε τη συνάρτηση ως προς τη μονοτονία και να αποδείξετε ότι $f(x) < 1$ για κάθε $x > 0$.
  \item Ένα σημείο $M$ κινείται επί της $C_f$ στο διάστημα $(0, +\infty)$. Αν η τετμημένη του αυξάνεται με σταθερό ρυθμό $2 \mu/\si{sec}$, να βρείτε το ρυθμό μεταβολής της τεταγμένης του $M$ τη χρονική στιγμή που διέρχεται από το σημείο με τετμημένη $x=1$.
\end{erwthma}
\end{thema}
%=========== ΤΕΛΟΣ - 22ο ΘΕΜΑ Γ ===========

%=========== 23ο ΘΕΜΑ Γ ===========
\begin{thema}{Γ}
Δίνεται η συνάρτηση $f(x) = \dfrac{e^x - e^{-x}}{2}$. 
\begin{erwthma}
  \item Να μελετήσετε την $f(x)$ ως προς τη μονοτονία, να βρείτε το σύνολο τιμών της και να δείξετε ότι αντιστρέφεται.
  \item Θεωρώντας γνωστό τον τύπο της αντίστροφης $f^{-1}(x) = \ln(x + \sqrt{x^2+1})$, να δείξετε ότι η εξίσωση $f(x) = 1$ έχει ακριβώς μία λύση.
  \item Να υπολογίσετε τα όρια στο $+\infty$ και $-\infty$ και να βεβαιώσετε αν έχει οριζόντιες ασύμπτωτες.
  \item Αφού υπολογίσετε το ολοκλήρωμα $\dint_0^1 f(x) \d x$, να βρείτε το εμβαδόν του χωρίου που περικλείεται από τη γραφική παράσταση της $f^{-1}$, τον άξονα $x'x$ και την κατακόρυφη ευθεία $x = \dfrac{e^2-1}{2e}$.
\end{erwthma}
\end{thema}
%=========== ΤΕΛΟΣ - 23ο ΘΕΜΑ Γ ===========

%=========== 24ο ΘΕΜΑ Γ ===========
\begin{thema}{Γ}
Δίνεται η συνάρτηση $f(x) = e^x$.
\begin{erwthma}
  \item Να βρείτε την εξίσωση της εφαπτομένης της γραφικής παράστασης της $f$ που διέρχεται από την αρχή των αξόνων $O(0,0)$.
  \item Μελετήστε την $f$ ως προς την κυρτότητα και αποδείξτε ότι $e^x \ge ex$ για κάθε πραγματικό αριθμό $x$.
  \item Να υπολογίσετε το όριο $\lim\limits_{x \to 1} \dfrac{\ln(f(x)) - 1}{x - 1}$.
  \item Να υπολογίσετε το όριο $\lim\limits_{x \to +\infty} \dfrac{f(x)}{f'(x) + x}$.
\end{erwthma}
\end{thema}
%=========== ΤΕΛΟΣ - 24ο ΘΕΜΑ Γ ===========

%=========== 25ο ΘΕΜΑ Γ ===========
\begin{thema}{Γ}
Μια συνεχής συνάρτηση $f: \mathbb{R} \to \mathbb{R}$ ικανοποιεί τη σχέση $f(x) = e^x + 2\dint_0^1 f(t)dt$ για κάθε $x \in \mathbb{R}$.
\begin{erwthma}
  \item Να αποδείξετε ότι ο τύπος της συνάρτησης είναι $f(x) = e^x + 2 - 2e$.
  \item Να βρείτε την εξίσωση της εφαπτομένης της $C_f$ στο σημείο με τετμημένη $x_0 = 1$ και να αποδείξετε ότι ισχύει $e^{x-1} \ge x$ για κάθε $x \in \mathbb{R}$.
  \item Να αποδείξετε ότι η εξίσωση $f(x) = x$ έχει ακριβώς μία ρίζα στο διάστημα $(1, 2)$.
  \item Να υπολογίσετε το εμβαδόν του χωρίου που περικλείεται από τη γραφική παράσταση της συνάρτησης $f$, την εφαπτομένη του ερωτήματος Γ2 και τον άξονα $y'y$.
\end{erwthma}
\end{thema}
%=========== ΤΕΛΟΣ - 25ο ΘΕΜΑ Γ ===========

%=========== 26ο ΘΕΜΑ Γ ===========
\begin{thema}{Γ}
Μια συνάρτηση $f$ είναι συνεχής στο $\mathbb{R}$ και ικανοποιεί την σχέση $|f(x)| \le x^2$ για κάθε $x \in \mathbb{R}$.
\begin{erwthma}
  \item Nα αποδείξετε ότι $f(0) = 0$.
  \item Να αποδείξετε, ότι η $f$ είναι παραγωγίσιμη στο $x_0=0$ με $f'(0) = 0$.
  \item Να υπολογίσετε το όριο $\lim\limits_{x \to 0} \dfrac{f(x) + \hm x}{x + \ef x}$.
  \item Αν η $f(x)$ είναι δύο φορές παραγωγίσιμη και κυρτή στο $(0, +\infty)$, να αποδείξετε ότι $f(x) \ge 0$ για $x > 0$.
\end{erwthma}
\end{thema}
%=========== ΤΕΛΟΣ - 26ο ΘΕΜΑ Γ ===========

%=========== 27ο ΘΕΜΑ Γ ===========
\begin{thema}{Γ}
Δίνεται η πολυωνυμική συνάρτηση $f(x) = \dfrac{1}{3}x^3 - \lambda x^2 + 4x - 2$, όπου $\lambda \in \mathbb{R}$.
\begin{erwthma}
  \item Να βρείτε το σύνολο των τιμών της παραμέτρου $\lambda$ ώστε η συνάρτηση να είναι γνησίως αύξουσα σε όλο το $\mathbb{R}$.
  \item Για $\lambda = 2$, να βρείτε το μοναδικό σημείο καμπής της γραφικής παράστασης της συνάρτησης.
  \item Για $\lambda = 2$, να αποδείξετε ότι $f(x) - \dfrac{2}{3} = \dfrac{1}{3}(x-2)^3$ και στη συνέχεια να λύσετε την ανίσωση $f(|x| - 1) < \dfrac{2}{3}$.
  \item Για $\lambda = 2$, να υπολογίσετε το εμβαδόν του χωρίου που περικλείεται από τη γραφική παράσταση της συνάρτησης $f$, την εφαπτομένη της στο σημείο καμπής και τον άξονα $y'y$.
\end{erwthma}
\end{thema}
%=========== ΤΕΛΟΣ - 27ο ΘΕΜΑ Γ ===========

%=========== 28ο ΘΕΜΑ Γ ===========
\begin{thema}{Γ}
Έστω η συνάρτηση $f(x) = \ln x$ με $x > 0$ και η ευθεία $(\varepsilon): y = x - 2$.
\begin{erwthma}
  \item Να βρείτε το σημείο $M$ της γραφικής παράστασης της $f$ στο οποίο η εφαπτομένη είναι παράλληλη προς την $\varepsilon$. 
  \item Βρείτε ποια είναι η ελάχιστη απόσταση της καμπύλης $C_f$ από την ευθεία $\varepsilon$.
  \item Ένα σημείο $N$ κινείται πάνω στην ευθεία $\varepsilon$ και η τετμημένη του αυξάνεται με ρυθμό $1\ \mu/\si{sec}$. Ποιος είναι ο ρυθμός μεταβολής της απόστασης $MN$ (όπου $M(1,0)$) τη χρονική στιγμή που το $N$ βρίσκεται στο σημείο με $x=2$?
  \item Να υπολογίσετε το όριο $\lim\limits_{x \to 1} \dfrac{f(\sqrt{x})}{x-1}$.
\end{erwthma}
\end{thema}
%=========== ΤΕΛΟΣ - 28ο ΘΕΜΑ Γ ===========

%=========== 29ο ΘΕΜΑ Γ ===========
\begin{thema}{Γ}
Για μια παραγωγίσιμη συνάρτηση $f: \mathbb{R} \to \mathbb{R}$ ισχύει $f'(x) - f(x) = e^x$ για κάθε πραγματικό αριθμό $x$, με αρχική συνθήκη $f(0) = 1$.
\begin{erwthma}
  \item Να αποδείξετε ότι $f(x) = (x+1)e^x$.
  \item Να βρείτε το μοναδικό ολικό ακρότατο της συνάρτησης $f(x)$.
  \item Να ελέγξετε τη συνάρτηση για οριζόντια ασύμπτωτη στο $-\infty$.
  \item Να υπολογίσετε το όριο: $\lim\limits_{x \to +\infty} \left[ f(x+1) - e \cdot f(x) \right]$.
\end{erwthma}
\end{thema}
%=========== ΤΕΛΟΣ - 29ο ΘΕΜΑ Γ ===========

%=========== 30ο ΘΕΜΑ Γ ===========
\begin{thema}{Γ}
Θεωρούμε τη συνάρτηση $f(x) = 2x^3 - 3a x^2 + 5$, όπου $a \ne 0$.
\begin{erwthma}
  \item Να βρείτε, συναρτήσει του $a$, τις θέσεις των δύο τοπικών ακροτάτων της $f(x)$.
  \item Να αποδείξετε ότι καθώς το $a$ μεταβάλλεται, το τοπικό ακρότατο που εξαρτάται από το $a$, δηλαδή το σημείο $M(a, f(a))$, κινείται πάνω στην καμπύλη με εξίσωση $y = 5 - x^3$.
  \item Να βρείτε την τιμή του $a$ για την οποία η γραφική παράσταση της $f$ παρουσιάζει σημείο καμπής στο σημείο με τετμημένη $x=1$.
  \item Για $a = 2$, να υπολογίσετε το όριο $\lim\limits_{x \to 0} \dfrac{f(x) - 5}{f'(x)}$.
\end{erwthma}
\end{thema}
%=========== ΤΕΛΟΣ - 30ο ΘΕΜΑ Γ ===========

%=========== 31ο ΘΕΜΑ Γ ===========
\begin{thema}{Γ}
Δίνεται η συνάρτηση $f(x) = \dfrac{x}{\sqrt{x^2+1}}$ για κάθε $x \in \mathbb{R}$.
\begin{erwthma}
  \item Να αποδείξετε ότι η συνάρτηση $f$ είναι γνησίως αύξουσα στο $\mathbb{R}$ και να βρείτε το σύνολο τιμών της.
  \item Να αποδείξετε ότι η συνάρτηση $f$ αντιστρέφεται και να βρείτε τον τύπο της αντίστροφης συνάρτησης $f^{-1}(x)$.
  \item Να βρείτε την εξίσωση της εφαπτομένης της γραφικής παράστασης της αντίστροφης συνάρτησης $f^{-1}$ στο σημείο της με τετμημένη $x_0 = \dfrac{\sqrt{2}}{2}$.
  \item Να υπολογίσετε το όριο $\lim\limits_{x \to 0} \dfrac{f^{-1}(x) - x}{x^3}$.
\end{erwthma}
\end{thema}
%=========== ΤΕΛΟΣ - 31ο ΘΕΜΑ Γ ===========

%=========== 32ο ΘΕΜΑ Γ ===========
\begin{thema}{Γ}
Έστω η συνάρτηση $f(x) = e^x - a x$ όπου $a \in \mathbb{R}$.
\begin{erwthma}
  \item Να υπολογίσετε την τιμή του $a$ ώστε να ισχύει $\lim\limits_{x \to 0} \dfrac{f(x) - 1}{x} = 2$.
  \item Για $a = 1$, να υπολογίσετε το όριο $\lim\limits_{x \to 0} \dfrac{f(x) - 1}{x^2}$.
  \item Για την παραπάνω συνάρτηση, αποδείξτε ότι $f(x) > 0$ για κάθε $x\in\mathbb{R}$.
  \item Να βρείτε το όριο $\lim\limits_{x \to +\infty} \left[ \dfrac{f(x)}{e^x} \cdot \ln x \right]$.
\end{erwthma}
\end{thema}
%=========== ΤΕΛΟΣ - 32ο ΘΕΜΑ Γ ===========

%=========== 33ο ΘΕΜΑ Γ ===========
\begin{thema}{Γ}
Δίνονται οι συναρτήσεις $f(x) = x^2 + 1$ και $g(x) = -x^2 + 4x - 1$.
\begin{erwthma}
  \item Να βρείτε το μοναδικό κοινό σημείο επαφής $A$ των γραφικών παραστάσεων $C_f$ και $C_g$.
  \item Να ελέγξετε αν οι δύο γραφικές παραστάσεις δέχονται κοινή εφαπτόμενη στο $A$ και αν ναι, να τη βρείτε.
  \item Να υπολογίσετε το όριο $\lim\limits_{x \to 1} \dfrac{\sqrt{f(x)} - \sqrt{2}}{x - 1}$.
  \item Να υπολογίσετε το εμβαδόν του χωρίου που περικλείεται από τις δύο καμπύλες και τον άξονα $y'y$.
\end{erwthma}
\end{thema}
%=========== ΤΕΛΟΣ - 33ο ΘΕΜΑ Γ ===========

%=========== 34ο ΘΕΜΑ Γ ===========
\begin{thema}{Γ}
Θεωρούμε τη συνάρτηση $f(x) = \begin{cdcases} x \ln x, & x > 0 \\ 0, & x = 0 \end{cdcases}$.
\begin{erwthma}
  \item Να υπολογίσετε το όριο $\lim\limits_{x\to 0^+} (x \ln x)$ και να εξετάσετε αν η $f$ είναι συνεχής στο $0$.
  \item Εξετάστε αν η συνάρτηση $f$ είναι παραγωγίσιμη στο $0$. Τι συμβαίνει με την εφαπτομένη της γραφικής παράστασης στο $(0,0)$?
  \item Υπολογίστε το όριο $\lim\limits_{x \to 0^+} [ f(x) \cdot f'(x) ]$.
  \item Υπολογίστε το ορισμένο ολοκλήρωμα $\dint_{\frac{1}{e}}^1 f(x)\d x$.
\end{erwthma}
\end{thema}
%=========== ΤΕΛΟΣ - 34ο ΘΕΜΑ Γ ===========

%=========== 35ο ΘΕΜΑ Γ ===========
\begin{thema}{Γ}
Έστω μια παραγωγίσιμη συνάρτηση $f: \mathbb{R} \to \mathbb{R}$ τέτοια ώστε για κάθε $x \in \mathbb{R}$ να ισχύει:
$$ e^{f(x)} + f(x) = x + 2 $$
\begin{erwthma}
  \item Να αποδείξετε ότι η συνάρτηση $f$ είναι γνησίως αύξουσα στο $\mathbb{R}$.
  \item Να αποδείξετε ότι η συνάρτηση $f$ αντιστρέφεται και να βρείτε το πεδίο ορισμού της $f^{-1}$.
  \item Να βρείτε τον τύπο της αντίστροφης συνάρτησης $f^{-1}(x)$.
  \item Να υπολογίσετε το εμβαδόν του χωρίου που περικλείεται από τη γραφική παράσταση της συνάρτησης $f$, τον άξονα $x'x$ και τις ευθείες με εξισώσεις $x = -1$ και $x = e - 1$.
\end{erwthma}
\end{thema}
%=========== ΤΕΛΟΣ - 35ο ΘΕΜΑ Γ ===========

%=========== 36ο ΘΕΜΑ Γ ===========
\begin{thema}{Γ}
Δίνεται η συνάρτηση $f(x) = \dfrac{e^x}{x^2+1}$ με πεδίο ορισμού το $\mathbb{R}$.
\begin{erwthma}
  \item Να μελετήσετε την $f$ ως προς τη μονοτονία και να βρείτε το σύνολο τιμών της.
  \item Να αποδείξετε ότι η εξίσωση $e^x = \lambda(x^2+1)$ έχει ακριβώς μία πραγματική ρίζα στο $\mathbb{R}$, για κάθε $\lambda > 0$.
  \item Να αποδείξετε ότι $1 < \dint_0^1 f(x) \d x < \dfrac{e}{2}$.
  \item Θεωρούμε συνάρτηση $F$, η οποία είναι μια παράγουσα της $f$ στο $\mathbb{R}$. Να αποδείξετε ότι υπάρχει ένα τουλάχιστον $x_0 \in (0, 1)$ τέτοιο, ώστε να ισχύει: $$F(x_0) + f(x_0) = F(1) + 1$$
\end{erwthma}
\end{thema}
%=========== ΤΕΛΟΣ - 36ο ΘΕΜΑ Γ ===========

%=========== 37ο ΘΕΜΑ Γ ===========
\begin{thema}{Γ}
Για κάθε $a \in \mathbb{R}$, θεωρούμε τη συνάρτηση $f(x) = \dfrac{\sqrt{x^2+1} - a x}{x}, x > 0$.
\begin{erwthma}
  \item Να βρείτε, συναρτήσει του $a$, το όριο $\lim\limits_{x \to +\infty} f(x)$.
  \item Να βρείτε τη μοναδική τιμή του $a$ ώστε η γραφική παράσταση της $f$ να έχει ως οριζόντια ασύμπτωτη τον άξονα $x'x$.
  \item Για $a = 1$, να υπολογίσετε το όριο $\lim\limits_{x \to +\infty} [ x f(x) ]$.
  \item Για $a=1$, να αποδείξετε ότι η συνάρτηση $f$ αντιστρέφεται και να βρείτε το πεδίο ορισμού της αντίστροφής της.
\end{erwthma}
\end{thema}
%=========== ΤΕΛΟΣ - 37ο ΘΕΜΑ Γ ===========

%=========== 38ο ΘΕΜΑ Γ ===========
\begin{thema}{Γ}
Έστω η συνάρτηση $f(x) = x^4 - 2x^2 + \lambda$, $\lambda > 0$. Γνωρίζουμε ότι ο άξονας $x'x$ εφάπτεται στη γραφική παράσταση της συνάρτησης ακριβώς στα τοπικά της ελάχιστα.
\begin{erwthma}
  \item Αποδείξτε ότι $\lambda = 1$.
  \item Να δείξετε ότι $\dint_{0}^{a}f(x) \ge 0$ για κάθε $a\in\mathbb{R}$.
  \item Βρείτε το πλήθος των ριζών της εξίσωσης $f(x) = k$ για $k \in (0, 1)$.
  \item Εξετάστε την ύπαρξη του ορίου $\lim\limits_{x \to 1} \dfrac{\sqrt{f(x)}}{x-1}$.
\end{erwthma}
\end{thema}
%=========== ΤΕΛΟΣ - 38ο ΘΕΜΑ Γ ===========

%=========== 39ο ΘΕΜΑ Γ ===========
\begin{thema}{Γ}
Για μια συνάρτηση $f: \mathbb{R} \to \mathbb{R}$ ισχύει η ανισότητα $|f(x) - f(y)| \le (x-y)^2$ για κάθε ζεύγος πραγματικών αριθμών $x, y$.
\begin{erwthma}
  \item Να δείξετε ότι η συνάρτηση είναι παραγωγίσιμη σε κάθε $x_0 \in \mathbb{R}$ με $f'(x_0) = 0$.
  \item Τι συμπεραίνετε για τον τύπο της $f(x)$?
  \item Αν γνωρίζουμε επίσης ότι $f(0) = 2$, να υπολογίσετε το όριο $\lim\limits_{x \to 0} \dfrac{\hm(f(x) - 2)}{x}$.
  \item Να υπολογίσετε το όριο $\lim\limits_{x \to +\infty} [ x (f(x+1) - f(x)) ]$.
\end{erwthma}
\end{thema}
%=========== ΤΕΛΟΣ - 39ο ΘΕΜΑ Γ ===========

%=========== 40ο ΘΕΜΑ Γ ===========
\begin{thema}{Γ}
Δίνεται η συνάρτηση $f(x) = x + \dfrac{1}{x}, x > 0$.
\begin{erwthma}
  \item Να βρείτε όλες τις ασύμπτωτες της γραφικής παράστασης (κατακόρυφες και πλάγιες).
  \item Έστω $a > 0$. Μια κατακόρυφη ευθεία $x = a$ τέμνει την $C_f$ στο σημείο $A$ και την πλάγια ασύμπτωτη στο σημείο $B$. Να βρείτε το μήκος του τμήματος $AB$ συναρτήσει του $a$.
  \item Να βρείτε το όριο $\lim\limits_{a \to +\infty} [a \cdot (AB)]$.
  \item Ένα σημείο $M$ κινείται στη γραφική παράσταση της $f$ έτσι ώστε η τετμημένη του να αυξάνεται με σταθερό ρυθμό $2\mu/\si{sec}$. Βρείτε το ρυθμό μεταβολής της τεταγμένης του $M$ τη χρονική στιγμή που το σημείο διέρχεται από τη θέση $(3,f(3))$.
\end{erwthma}
\end{thema}
%=========== ΤΕΛΟΣ - 40ο ΘΕΜΑ Γ ===========

%=========== 41ο ΘΕΜΑ Γ ===========
\begin{thema}{Γ}
Έστω η συνάρτηση $f(x) = \ln(x^2 + 1)$.
\begin{erwthma}
  \item Να μελετήσετε τη συνάρτηση $f$ ως προς τη μονοτονία, τα ακρότατα, την κυρτότητα και τα σημεία καμπής.
  \item Να αποδείξετε ότι $f(x) \le x^2$ για κάθε $x \in \mathbb{R}$.
  \item Να υπολογίσετε το όριο $\lim\limits_{x \to +\infty} \dfrac{f(x)}{x}$.
  \item Να αποδείξετε ότι η συνάρτηση $f'$ είναι περιττή και να υπολογίσετε το εμβαδόν του χωρίου που περικλείεται από τη γραφική παράσταση της $f'$, τον άξονα $x'x$ και τις κατακόρυφες ευθείες $x=-1$ και $x=1$.
\end{erwthma}
\end{thema}
%=========== ΤΕΛΟΣ - 41ο ΘΕΜΑ Γ ===========

%=========== 42ο ΘΕΜΑ Γ ===========
\begin{thema}{Γ}
Μια συνάρτηση $f(x)$ είναι δύο φορές παραγωγίσιμη στο $[0, 1]$ με $f(0)=0$ και $f(1)=1$. Γνωρίζουμε επίσης ότι η εφαπτομένη της $C_f$ στην αρχή των αξόνων είναι ο άξονας $x'x$.
\begin{erwthma}
  \item Να εξηγήσετε γιατί $f'(0) = 0$.
  \item Να αποδείξετε ότι υπάρχει τουλάχιστον ένα $\xi \in (0, 1)$ ώστε $f'(\xi) = 1$.
  \item Αποδείξτε ότι η δεύτερη παράγωγος μηδενίζεται σε τουλάχιστον ένα εσωτερικό σημείο.
  \item Έστω ότι $f(x) < x$ για κάθε $x \in (0, 1)$. Να υπολογίσετε το $\lim\limits_{x \to 0^+} \dfrac{f(x) \ln x}{x}$.
\end{erwthma}
\end{thema}
%=========== ΤΕΛΟΣ - 42ο ΘΕΜΑ Γ ===========

%=========== 43ο ΘΕΜΑ Γ ===========
\begin{thema}{Γ}
Δίνεται η συνεχής συνάρτηση $f: \mathbb{R} \to \mathbb{R}$ για την οποία ισχύει η σχέση:
$$ f(x) = e^x - \dint_0^1 x f(t) \d t, \quad \text{για κάθε } x \in \mathbb{R} $$
\begin{erwthma}
  \item Να αποδείξετε ότι ο τύπος της συνάρτησης είναι $f(x) = e^x - \dfrac{2(e-1)}{3}x$.
  \item Να μελετήσετε την $f$ ως προς τη μονοτονία και να βρείτε τη θέση και το είδος του ακροτάτου.
  \item Να αποδείξετε ότι η εξίσωση $f(x) = 1$ έχει ακριβώς μία ρίζα στο ανοικτό διάστημα $(0, 1)$.
  \item Να αποδείξετε ότι υπάρχει ένα τουλάχιστον $\xi \in (0, 1)$ τέτοιο, ώστε η εφαπτομένη της γραφικής παράστασης της $f$ στο σημείο $M(\xi, f(\xi))$ να είναι παράλληλη στην ευθεία $\varepsilon: (1-e)x + 3y - 2026 = 0$.
\end{erwthma}
\end{thema}
%=========== ΤΕΛΟΣ - 43ο ΘΕΜΑ Γ ===========

%=========== 44ο ΘΕΜΑ Γ ===========
\begin{thema}{Γ}
Ένας κώνος είναι εγγεγραμμένος σε μια σφαίρα σταθερής ακτίνας $R=3$. Έστω $x$ η απόσταση του κέντρου της σφαίρας από τη βάση του κώνου, με $x \in (0, 3)$. 
\begin{center}
\begin{tikzpicture}
    % ==========================================
    % Ορισμός Παραμέτρων 
    % ==========================================
    \def\R{3}       % Ακτίνα σφαίρας R=3
    \def\d{1}       % Απόσταση x 
    \def\r{2.828}   % Ακτίνα βάσης
    \def\yell{0.25} % Προοπτική ελλείψεων 

    % ==========================================
    % 1. Σχεδιασμός Σφαίρας (με Radial Shading)
    % ==========================================
    \shadedraw[shading=radial, inner color=white, outer color=gray!30, draw=gray!80, thick] (0,0) circle (\R);
    
    % Προαιρετικά: Προσθέτουμε ένα "γυάλισμα" 
    \fill[white, opacity=0.3] (-0.3*\R, 0.4*\R) ellipse ({0.2*\R} and {0.15*\R});

    % Ισημερινός σφαίρας 
    \draw[dashed, color=gray!60] (\R,0) arc (0:180:{\R} and {\R*\yell});
    \draw[color=gray!60] (-\R,0) arc (180:360:{\R} and {\R*\yell});

    % ==========================================
    % 2. Σχεδιασμός Βάσης Κώνου 
    % ==========================================
    \coordinate (K) at (0, -\d);
    
    \shadedraw[dashed, thick, color=\xrwma!60!black, 
               top color=\xrwma!10, bottom color=\xrwma!20, fill opacity=0.6] 
        (\r, -\d) arc (0:180:{\r} and {\r*\yell})
        arc (180:360:{\r} and {\r*\yell});

    % ==========================================
    % 3. Σχεδιασμός Πλευρών Κώνου (με Linear Shading)
    % ==========================================
    \coordinate (V) at (0, \R);               
    \coordinate (Left) at (-\r, -\d);         
    \coordinate (Right) at (\r, -\d);         
    
    % Shading πλευρών
    \shadedraw[thick, color=\xrwma!60!black,
               left color=\xrwma!60, right color=\xrwma!60, middle color=white!80!\xrwma, fill opacity=0.8] 
               (V) -- (Left) arc (180:360:{\r} and {-\r*\yell}) -- cycle;

    % ==========================================
    % 4. Γεωμετρικά Στοιχεία & Ετικέτες
    % ==========================================
    \coordinate (O) at (0, 0);                
    \coordinate (A) at (\r, -\d);             
    
    % Κατακόρυφος άξονας
    \draw[dashed, thick] (V) -- (O);
    \draw[thick, secondcolor] (O) -- (K) node[midway, left] {$x$};
    
    % Ακτίνα βάσης κώνου (ρ)
    \draw[dashed, thick, \xrwma!60!black] (K) -- (A) node[midway, below] {$\rho$};
    
    % Ακτίνα σφαίρας (R)
    \draw[thick, \xrwma!50!black] (O) -- (A) node[midway, above] {$R=3$};

    % Σημεία
    \filldraw (O) circle (1pt) node[left] {$O$};
    \filldraw (K) circle (1pt) node[below left] {$K$};
    \filldraw (V) circle (1pt) node[above] {$V$};
    \filldraw (A) circle (1pt) node[right] {$A$};

\end{tikzpicture}
\end{center}
\begin{erwthma}
  \item Να αποδείξετε ότι η ακτίνα $\rho$ της βάσης του κώνου δίνεται από τη σχέση $\rho^2 = 9 - x^2$.
  \item Αποδείξτε ότι ο όγκος του κώνου συναρτήσει του $x$ δίνεται από τον τύπο $V(x) = \dfrac{\pi}{3}(9-x^2)(3+x)$.
  \item Να βρείτε την τιμή του $x$ για την οποία ο όγκος του κώνου μεγιστοποιείται, καθώς και τις διαστάσεις (ύψος και ακτίνα) του μέγιστου κώνου.
  \item Αν το $x$ αυξάνεται με ρυθμό $0.5 \si{cm/s}$, να βρείτε το ρυθμό μεταβολής του όγκου $V(x)$ τη χρονική στιγμή που το ύψος του κώνου είναι $4\si{cm}$.
\end{erwthma}
\end{thema}
%=========== ΤΕΛΟΣ - 44ο ΘΕΜΑ Γ ===========

%=========== 45ο ΘΕΜΑ Γ ===========
\begin{thema}{Γ}
Θεωρούμε τη συνάρτηση $f: \mathbb{R} \to \mathbb{R}$ η οποία για κάθε $x \in \mathbb{R}$ ικανοποιεί τη σχέση \[f^3(x) + 3f^2(x)\hm x + 3f(x)\hm^2 x = x^3 - \hm^3 x.\]
\begin{erwthma}
  \item Να αποδείξετε ότι $f(x) = x - \hm x$.
  \item Να αποδείξετε ότι για κάθε $x > 0$ ισχύει $f(x) > 0$.
  \item Να μελετήσετε την $f(x)$ ως προς την κυρτότητα στο διάστημα $[0, \pi]$.
  \item Να υπολογίσετε το εμβαδόν του χωρίου που περικλείεται από τη γραφική παράσταση της $f$, τον άξονα $x'x$ και την ευθεία $x = \frac{\pi}{2}$.
\end{erwthma}
\end{thema}
%=========== ΤΕΛΟΣ - 45ο ΘΕΜΑ Γ ===========

%=========== 46ο ΘΕΜΑ Γ ===========
\begin{thema}{Γ}
Δίνεται η συνάρτηση $f(x) = e^{x-1} + x^2 - 2$ για κάθε $x \in \mathbb{R}$.
\begin{erwthma}
  \item Να αποδείξετε ότι η $f$ είναι κυρτή στο $\mathbb{R}$.
  \item Να βρείτε την εξίσωση της εφαπτομένης της $C_f$ στο σημείο της με τετμημένη $x_0 = 1$.
  \item Να υπολογίσετε το όριο $\lim\limits_{x \to 1} \dfrac{f(x) - 3x + 3}{\ln^2 x}$.
  \item Να αποδείξετε ότι η εξίσωση $f(x) + x f'(x) = 0$ έχει τουλάχιστον μία ρίζα στο ανοικτό διάστημα $(0, 1)$.
\end{erwthma}
\end{thema}
%=========== ΤΕΛΟΣ - 46ο ΘΕΜΑ Γ ===========

%=========== 47ο ΘΕΜΑ Γ ===========
\begin{thema}{Γ}
Σε έναν φανοστάτη ύψους $H = 6$ \si{m} πλησιάζει ένας άνθρωπος ύψους $h = 2$ \si{m} με σταθερή ταχύτητα $u = 1.5$ \si{m/s} επί ευθύγραμμου οριζόντιου δρόμου. 
\begin{center}
\begin{tikzpicture}[scale=1.2, font=\sffamily]

    % ==========================================
    % Παράμετροι Σχήματος
    % ==========================================
    \def\H{6.0}   % Ύψος φανοστάτη
    \def\h{2.0}   % Ύψος ανθρώπου
    % Για το σχήμα επιλέγουμε μια απόσταση x=4. 
    % Με όμοια τρίγωνα το μήκος της σκιάς s είναι: s/2 = (4+s)/6 => 6s = 8 + 2s => 4s = 8 => s = 2
    \def\x{4.0}   % Απόσταση ανθρώπου από φανοστάτη
    \def\s{2.0}   % Μήκος σκιάς

    % ==========================================
    % Συντεταγμένες
    % ==========================================
    \coordinate (L_base) at (0, 0);                 % Βάση Φανοστάτη
    \coordinate (L_top)  at (0, \H);                % Λάμπα
    
    \coordinate (M_base) at (\x, 0);                % Βάση Ανθρώπου
    \coordinate (M_top)  at (\x, \h);               % Κεφάλι Ανθρώπου
    
    \coordinate (S_tip)  at (\x + \s, 0);           % Άκρη Σκιάς (x+s)

    % ==========================================
    % 1. Σχεδίαση Εδάφους & Αντικειμένων
    % ==========================================
    % Έδαφος
    \draw[thick, fill=gray!20] (-1, -0.3) rectangle (\x + \s + 2, 0);
    \draw[very thick] (-1, 0) -- (\x + \s + 2, 0);

    % Φανοστάτης (Στύλος και Λάμπα)
    \draw[line width=4pt, gray!80!black] (L_base) -- (L_top);
    \fill[secondcolor!80!secondcolor] (L_top) circle (0.15);
    \filldraw[fill=secondcolor!20, draw=secondcolor!50!secondcolor, thick, opacity=0.7] (L_top) circle (0.4); % Λάμψη

    % Άνθρωπος (Αναπαράσταση με χοντρή γραμμή)
    \draw[line width=3pt, \xrwma!70!black] (M_base) -- (M_top);
    \filldraw[\xrwma!70!black] (\x, \h-0.15) circle (0.15); % Κεφάλι

    % Σκιά (Ημι-διαφανής μαύρη γραμμή στο έδαφος)
    \draw[line width=4pt, black, opacity=0.3] (M_base) -- (S_tip);

    % ==========================================
    % 2. Φωτεινή Ακτίνα
    % ==========================================
    \draw[dashed, thick, secondcolor] (L_top) -- (M_top) -- (S_tip);

    % ==========================================
    % 3. Ετικέτες, Διαστάσεις και Ταχύτητα
    % ==========================================
    
    % Ταχύτητα (Διάνυσμα υ)
    \draw[->, >=latex, ultra thick, secondcolor] (\x, \h/2) -- (\x-1.5, \h/2) node[above] {$u = 1.5$ \si{m/s}};

    % Ύψος Φανοστάτη (H)
    \draw[<->, >=latex] (-0.6, 0) -- (-0.6, \H) node[midway, left] {$H = 6$ \si{m}};
    \draw[dashed, gray] (L_base) -- (-0.8, 0);
    \draw[dashed, gray] (L_top) -- (-0.8, \H);

    % Ύψος Ανθρώπου (h)
    \draw[<->, >=latex] (\x + 0.5, 0) -- (\x + 0.5, \h) node[midway, right] {$h = 2$ \si{m}};

    % Απόσταση (x)
    \draw[<->, >=latex] (0, -0.6) -- (\x, -0.6) node[midway, below] {$x$};
    \draw[dashed, gray] (L_base) -- (0, -0.8);
    \draw[dashed, gray] (M_base) -- (\x, -0.8);

    % Μήκος Σκιάς (s)
    \draw[<->, >=latex] (\x, -0.6) -- (\x + \s, -0.6) node[midway, below] {$s$};
    \draw[dashed, gray] (S_tip) -- (\x + \s, -0.8);

    % Συνολική Απόσταση (y = x + s)
    \draw[<->, >=latex, \xrwma!70!black] (0, -1.2) -- (\x + \s, -1.2) node[midway, below] {$y = x + s$};
    \draw[dashed, \xrwma!70!black] (0, -0.8) -- (0, -1.4);
    \draw[dashed, \xrwma!70!black] (\x + \s, -0.8) -- (\x + \s, -1.4);

    % Κουκκίδες Σημείων
    \filldraw[black] (L_base) circle (1.5pt) node[above right] {$O$};
    \filldraw[black] (M_base) circle (1.5pt) node[above right] {$A$};
    \filldraw[black] (S_tip) circle (1.5pt) node[above right] {$\Gamma$};

\end{tikzpicture}
\end{center}
\begin{erwthma}
  \item Αν τη χρονική στιγμή $t$ ο άνθρωπος απέχει απόσταση $x(t)$ από τον φανοστάτη, να εκφράσετε το μήκος της σκιάς του $s(t)$ συναρτήσει του $x(t)$.
  \item Αποδείξτε ότι η άκρη της σκιάς του ανθρώπου κινείται με σταθερή ταχύτητα ως προς τη βάση του φανοστάτη. 
  \item Έστω $\theta(t)$ η γωνία που σχηματίζει η οπτική ακτίνα από την κορυφή του φανοστάτη στο κεφάλι του ανθρώπου με την κορυφή του φανοστάτη. Εκφράστε την $\ef\theta$ συναρτήσει του $x$.
  \item Βρείτε το ρυθμό μεταβολής της γωνίας $\theta$ τη χρονική στιγμή που ο άνθρωπος απέχει $4$ m από τον φανοστάτη.
\end{erwthma}
\end{thema}
%=========== ΤΕΛΟΣ - 47ο ΘΕΜΑ Γ ===========

%=========== 48ο ΘΕΜΑ Γ ===========
\begin{thema}{Γ}
Γνωρίζουμε για τη συνάρτηση $f(x)$ ότι είναι δύο φορές παραγωγίσιμη στο $\mathbb{R}$ με $f(0) = 0$ και $f'(0) = 0$. Η γραφική παράσταση της δεύτερης παραγώγου $f''(x)$ αποτελείται από δύο ημιευθείες: από το $-\infty$ μέχρι το $x=2$ είναι $y = 1$, ενώ από το $x=2$ έως το $+\infty$ ταυτίζεται με τον άξονα $x'x$.
\begin{erwthma}
  \item Να γράψετε τον τύπο της δεύτερης παραγώγου $f''(x)$.
  \item Γνωρίζοντας ότι $f'(0) = 0$, να βρείτε τον τύπο της $f'(x)$. Πού παρουσιάζει ελάχιστο?
  \item Να βρείτε τον τύπο της $f(x)$. Είναι συνεχής στο $x=2$?
  \item Να μελετήσετε την $f(x)$ ως προς τη μονοτονία και τα τοπικά ακρότατα, αν υπάρχουν.
\end{erwthma}
\end{thema}
%=========== ΤΕΛΟΣ - 48ο ΘΕΜΑ Γ ===========

%=========== 49ο ΘΕΜΑ Γ ===========
\begin{thema}{Γ}
Δίνεται η παραβολή $C_f: f(x) = x^2 + a x + \beta, a, \beta \in \mathbb{R}$. Φέρνουμε δύο εφαπτόμενες $\varepsilon_1,\varepsilon_2$ την $C_f$ από το εξωτερικό σημείο $A(0, -1)$, οι οποίες εφάπτονται της $f$ στα σημεία $M(x_1, f(x_1))$ και $N(x_2, f(x_2))$.
\begin{erwthma}
  \item Να γράψετε την εξίσωση της εφαπτομένης της $C_f$ που διέρχεται από το $A(0,-1)$, ως προς το $x_0$ όπου $M(x_0,f(x_0))$ το σημείο επαφής.
  \item Αν γνωρίζουμε ότι οι δύο εφαπτόμενες είναι κάθετες μεταξύ τους, αποδείξτε ότι $\beta = -\dfrac{3}{4}$.
  \item Για $\beta = -3/4$ και $a = 0$, υπολογίστε το εμβαδόν μεταξύ των εφαπτομένων και της $C_f$.
\end{erwthma}
\end{thema}
%=========== ΤΕΛΟΣ - 49ο ΘΕΜΑ Γ ===========

%=========== 50ο ΘΕΜΑ Γ ===========
\begin{thema}{Γ}
Μια συνεχής συνάρτηση $f(x)$ ορίζεται στο $\mathbb{R}$ και $f(0)=2$. Γνωρίζουμε ότι η παράγωγός της ικανοποιεί τη σχέση $f'(x) = \dfrac{|x|}{x} e^{|x|}$ για κάθε $x \ne 0$.
\begin{erwthma}
  \item Να δείξετε ότι ο τύπος της $f$ είναι $f(x)=e^{|x|}+1$.
  \item Εξετάστε αν η συνάρτηση είναι παραγωγίσιμη στο σημείο $x_0=0$.
  \item Να μελετήσετε τη μονοτονία και τα ακρότατα της συναρτήσεως.
  \item Υπολογίστε το όριο $\lim\limits_{x \to -\infty} x^2 \cdot f(x)$.
\end{erwthma}
\end{thema}
%=========== ΤΕΛΟΣ - 50ο ΘΕΜΑ Γ ===========

%=========== 51ο ΘΕΜΑ Γ ===========
\begin{thema}{Γ}
Θεωρούμε τη συνάρτηση $h(x) = x^2 \ln\left(1 + \dfrac{1}{x}\right) - x$ με $x > 0$.
\begin{erwthma}
  \item Να υπολογίσετε τα όρια $\lim\limits_{x \to 0^+} h(x)$ και $\lim\limits_{x \to +\infty} h(x)$.
  \item Να αποδείξετε ότι η $h$ είναι κυρτή στο $(0, +\infty)$.
  \item Να αποδείξετε ότι η συνάρτηση $h$ είναι γνησίως φθίνουσα και να βρείτε το σύνολο τιμών της.
  \item Να αποδείξετε ότι $-\dfrac{1}{2} < \dint_1^2 h(x) \, \d x < 0$.
\end{erwthma}
\end{thema}
%=========== ΤΕΛΟΣ - 51ο ΘΕΜΑ Γ ===========

%=========== 52ο ΘΕΜΑ Γ ===========
\begin{thema}{Γ}
Έστω η γνησίως αύξουσα συνάρτηση $f: \mathbb{R} \to \mathbb{R}$ τέτοια ώστε $f(1) = 2$ και $f(2) = 3$. 
\begin{erwthma}
  \item Να εξηγήσετε γιατί η $f$ αντιστρέφεται και να βρείτε τις αντίστοιχες τιμές $f^{-1}(2)$ και $f^{-1}(3)$.
  \item Να λύσετε την εξίσωση $f^{-1}\big(f(x^2)-1\big) = 1$.
  \item Να λύσετε την ανίσωση $f\big(f^{-1}(2x) - 1\big) > 3$.
  \item Αν η γραφική παράσταση της $f$ διέρχεται από το σημείο $(0,0)$, να υπολογίσετε το ολοκλήρωμα $\dint_0^3 f^{-1}(y)dy$ εάν γνωρίζετε ότι $\dint_0^2 f(x)\d x = 4$.
\end{erwthma}
\end{thema}
%=========== ΤΕΛΟΣ - 52ο ΘΕΜΑ Γ ===========

%=========== 53ο ΘΕΜΑ Γ ===========
\begin{thema}{Γ}
Έχουμε ένα σύρμα συνολικού μήκους $10$ μέτρων. Κόβουμε το σύρμα σε δύο τμήματα ώστε με το ένα να κατασκευάσουμε ένα τετράγωνο και με το άλλο ένα ισόπλευρο τρίγωνο. Έστω $x$ η περίμετρος του τετραγώνου.
\begin{erwthma}
  \item Να εκφράσετε την πλευρά του τετραγώνου και την πλευρά του ισόπλευρου τριγώνου συναρτήσει του $x$.
  \item Αποδείξτε ότι το συνολικό εμβαδόν των δύο σχημάτων δίνεται από τη συνάρτηση $E(x) = \dfrac{x^2}{16} + \dfrac{\sqrt{3}(10-x)^2}{36}$.
  \item Για ποια τιμή του $x$ επιτυγχάνεται το ελάχιστο δυνατό συνολικό εμβαδόν? Υπολογίστε τον λόγο της πλευράς του τετραγώνου προς την πλευρά του τριγώνου σε αυτήν περίπτωση.
  \item Πως θα κόψουμε το σύρμα για να αποκτήσουμε το μέγιστο δυνατό εμβαδόν?
\end{erwthma}
\end{thema}
%=========== ΤΕΛΟΣ - 53ο ΘΕΜΑ Γ ===========

%=========== 54ο ΘΕΜΑ Γ ===========
\begin{thema}{Γ}
Έστω η παραγωγίσιμη συνάρτηση $f: [0, 2] \to \mathbb{R}$ με $f(0)=0, f(1)=1$ και $f(2)=2$.
\begin{erwthma}
  \item Να αποδείξετε ότι υπάρχουν $\xi_1 \in (0, 1)$ και $\xi_2 \in (1, 2)$ τέτοια ώστε $f'(\xi_1) = 1$ και $f'(\xi_2) = 1$.
  \item Αν η συνάρτηση είναι δύο φορές παραγωγίσιμη, δείξτε ότι έχει σημείο καμπής.
  \item Υποθέτοντας ότι $f''(x) > 0$ για $x \in (0,1)$ και $f''(x) < 0$ για $x \in (1,2)$ βρείτε σε ποιο σημείο μεγιστοποιείται ο ρυθμός μεταβολής της $f$.
  \item Βρείτε την εξίσωση της εφαπτομένης της $C_f$ στο σημείο $A(1,f(1))$ και αποδείξτε ότι $f(x) > x$ για κάθε $x \in (0, 1)$ και $f(x) < x$ για κάθε $x \in (1, 2)$.
\end{erwthma}
\end{thema}
%=========== ΤΕΛΟΣ - 54ο ΘΕΜΑ Γ ===========

%=========== 55ο ΘΕΜΑ Γ ===========
\begin{thema}{Γ}
Θεωρούμε την εξίσωση $e^x + x = 2$.
\begin{erwthma}
  \item Δείξτε ότι υπάρχει μια ρίζα $x_0$ της εξίσωσης στο διάστημα $(0, 1)$.
  \item Αποδείξτε ότι αυτή η ρίζα είναι μοναδική σε όλο το $\mathbb{R}$.
  \item Να μελετήσετε το πρόσημο της συνάρτησης $f(x)=e^x+x-2$.
  \item Θεωρούμε τη συνάρτηση $F(x) = e^x + \dfrac{x^2}{2} - 2x$, $x \in \mathbb{R}$. Να αποδείξετε ότι η $F$ παρουσιάζει ολικό ελάχιστο στη θέση $x_0$ και, στη συνέχεια, να υπολογίσετε το ολοκλήρωμα $\dint_0^{x_0} f(x) \d x$ ως πολυώνυμο του $x_0$.
\end{erwthma}
\end{thema}
%=========== ΤΕΛΟΣ - 55ο ΘΕΜΑ Γ ===========

%=========== 56ο ΘΕΜΑ Γ ===========
\begin{thema}{Γ}
Δίνεται η συνάρτηση $f(x) = \ln x$, ορισμένη στο $(0, +\infty)$.
\begin{erwthma}
  \item Να αποδείξετε ότι υπάρχει ακριβώς μία εφαπτομένη της γραφικής παράστασης $C_f$ η οποία διέρχεται από την αρχή των αξόνων $O(0, 0)$, και να βρείτε την εξίσωσή της.
  \item Να αποδείξετε ότι $\ln x \le \dfrac{x}{e}$ για κάθε $x > 0$ και στη συνέχεια να λύσετε στο $(0, +\infty)$ την εξίσωση $e^x = x^e$.
  \item Να βρείτε το σημείο της $C_f$ που απέχει την ελάχιστη απόσταση από την ευθεία $\zeta: y = x$ και να υπολογίσετε την απόσταση αυτή.
  \item Να υπολογίσετε το εμβαδόν του χωρίου που περικλείεται από την $C_f$, την εφαπτομένη του πρώτου ερωτήματος και τον άξονα $x'x$.
  \item Να υπολογίσετε το όριο $\lim\limits_{x \to 0^+} \left[ x f(x) \cdot \hm\left(\dfrac{1}{x}\right) \right]$.
\end{erwthma}
\end{thema}
%=========== ΤΕΛΟΣ - 56ο ΘΕΜΑ Γ ===========

%=========== 57ο ΘΕΜΑ Γ ===========
\begin{thema}{Γ}
Δίνεται η συνάρτηση $f(x) = x^3 + 2x - 3$, με $x \in \mathbb{R}$.
\begin{erwthma}
  \item Να αποδείξετε ότι η συνάρτηση $f$ αντιστρέφεται και να βρείτε το πεδίο ορισμού της αντίστροφης συνάρτησης $f^{-1}$.
  \item Να λύσετε στο $\mathbb{R}$ την ανίσωση $f\left(f^{-1}(x^2 - 4x) + 1\right) < 0$.
  \item Ένα υλικό σημείο $M(x, y)$ κινείται πάνω στη γραφική παράσταση της $f$. Η τετμημένη του αυξάνεται με σταθερό ρυθμό $2$ μονάδες/$\sec$. Τη χρονική στιγμή $t_0$ που το σημείο $M$ διέρχεται από το σημείο $M_0(2, 9)$, να βρείτε το ρυθμό μεταβολής του εμβαδού του τριγώνου $OAM$, όπου $O(0,0)$ η αρχή των αξόνων και $A(3, 0)$.
  \item Να υπολογίσετε το εμβαδόν του χωρίου που περικλείεται από τη γραφική παράσταση της $f^{-1}$, τους άξονες $x'x, y'y$ και την κατακόρυφη ευθεία $x = 9$.
\end{erwthma}
\end{thema}
%=========== ΤΕΛΟΣ - 57ο ΘΕΜΑ Γ ===========

%=========== 58ο ΘΕΜΑ Γ ===========
\begin{thema}{Γ}
Δίνεται το χωρίο $\Omega$ που ορίζεται από την παραβολή $y = x^2$ και τον άξονα $x'x$ με $x \in [0, 4]$.
\begin{erwthma}
  \item Να υπολογίσετε το συνολικό εμβαδόν του χωρίου $\Omega$.
  \item Η ευθεία $y = a x$ (με $a > 0$) διαιρεί το παραπάνω χωρίο σε δύο ισεμβαδικά μέρη $E_1, E_2$. Να εκφράσετε το εμβαδόν του πάνω μέρους συναρτήσει του $a$.
  \item Να βρείτε την τιμή της κλίσης $a$ που διχοτομεί το ολικό εμβαδόν $E$. 
  \item Να βρείτε το ρυθμό μεταβολής του εμβαδού του άνω χωρίου $E(a)$ αν η κλίση $a$ αυξάνεται με $1\ \mu/\sec$ τη στιγμή που $a=2$.
\end{erwthma}
\end{thema}
%=========== ΤΕΛΟΣ - 58ο ΘΕΜΑ Γ ===========

%=========== 59ο ΘΕΜΑ Γ ===========
\begin{thema}{Γ}
Έστω η εξίσωση $f(x) e^{f(x)} = x, x \ge 0, f(x) \ge 0$.
\begin{erwthma}
  \item Δείξτε ότι η συνάρτηση $f(x)$ είναι γνησίως αύξουσα στο διάστημα $[0, +\infty)$.
  \item Υπολογίστε την τιμή $f(0)$ και την τιμή $f(e)$. 
  \item Να βρείτε την εξίσωση της εφαπτομένης της $C_f$ στο σημείο με τετμημένη $x_0 = e$.
  \item Να αποδείξετε ότι $\lim\limits_{x \to +\infty} f(x) = +\infty$ και στη συνέχεια να υπολογίσετε το όριο $\lim\limits_{x \to +\infty} \dfrac{f(x)}{\ln x}$.
\end{erwthma}
\end{thema}
%=========== ΤΕΛΟΣ - 59ο ΘΕΜΑ Γ ===========

%=========== 60ο ΘΕΜΑ Γ ===========
\begin{thema}{Γ}
Δίνονται οι συναρτήσεις $f(x) = a \ln x$ και $g(x) = e^x - \beta$, με $x > 0$. Έστω ότι έχουν κοινή εφαπτομένη σε κοινό σημείο με $x_0 = 1$.
\begin{erwthma}
  \item Να αποδείξετε ότι $a = e$ και $\beta = e$.
  \item Να βρείτε την εξίσωση της κοινής εφαπτομένης $\varepsilon$ των γραφικών παραστάσεων των $f$ και $g$ στο σημείο με τετμημένη $x_0 = 1$.
  \item Να μελετήσετε τις $f$ και $g$ ως προς την κυρτότητα και να αποδείξετε ότι $f(x) \le e x - e \le g(x)$ για κάθε $x > 0$.
  \item Να αποδείξετε ότι $e^x - e \ln x \ge e$ για κάθε $x > 0$ και να υπολογίσετε το όριο $\lim\limits_{x \to 1} \dfrac{e^x - e\ln x - e}{(x-1)^2}$.
\end{erwthma}
\end{thema}
%=========== ΤΕΛΟΣ - 60ο ΘΕΜΑ Γ ===========

%=========== 61ο ΘΕΜΑ Γ ===========
\begin{thema}{Γ}
Ένα κινητό σημείο $M$ βρίσκεται αρχικά στην αρχή των αξόνων $O(0,0)$ και αρχίζει να κινείται κατά μήκος του θετικού ημιάξονα $Ox$ με σταθερή ταχύτητα $v_M = 3\si{m/s}$. Ταυτόχρονα, ένα δεύτερο σημείο $N$ ξεκινά να κινείται στον θετικό ημιάξονα $Oy$ ξεκινώντας από το $y_0 = 10$, επίσης με σταθερή ταχύτητα, απομακρυνόμενο προς το $+\infty$ ($v_N = 4\si{m/s}$). 
\begin{erwthma}
  \item Να εκφράσετε τις συντεταγμένες των σημείων $M(t), N(t)$ ως συναρτήσεις του χρόνου $t$.
  \item Εκφράστε την απόσταση $d(t) = (MN)$ μεταξύ των κινητών σημείων συναρτήσει του χρόνου $t$.
  \item Ποιος είναι ο ρυθμός μεταβολής της απόστασης $d(t)$ όταν $t = 2$ δευτερόλεπτα?
  \item Υπάρχει χρονική στιγμή κατά την οποία η απόσταση των δύο σημείων να γίνεται ελάχιστη?
\end{erwthma}
\end{thema}
%=========== ΤΕΛΟΣ - 61ο ΘΕΜΑ Γ ===========

%=========== 62ο ΘΕΜΑ Γ ===========
\begin{thema}{Γ}
Για μια συνάρτηση $f: \mathbb{R} \to \mathbb{R}$ δίνεται ότι ικανοποιεί την ανισότητα $f(x) - f(y) \ge e^x - e^y$ για κάθε $x \ge y$. 
\begin{erwthma}
  \item Να αποδείξετε ότι η συνάρτηση $g(x) = f(x) - e^x$ είναι αύξουσα στο $\mathbb{R}$.
  \item Αν επίσης γνωρίζουμε ότι η $f$ είναι παραγωγίσιμη, να αποδείξετε ότι $f'(x) \ge e^x$.
  \item Δείξτε ότι $f(x) - f(0) \ge e^x - 1$.
  \item Αν $f(0) = 1$, υπολογίστε το όριο $\lim\limits_{x \to +\infty} \dfrac{f(x)}{x}$.
\end{erwthma}
\end{thema}
%=========== ΤΕΛΟΣ - 62ο ΘΕΜΑ Γ ===========

%=========== 63ο ΘΕΜΑ Γ ===========
\begin{thema}{Γ}
Η συνάρτηση $f(x) = \dfrac{x^2+x+1}{x}$ ορίζεται στο $(0, +\infty)$.
\begin{erwthma}
  \item Να μελετήσετε τη συνάρτηση $f$ ως προς τη μονοτονία και τα τοπικά ακρότατα.
  \item Να βρείτε τις ασύμπτωτες της γραφικής παράστασης της $f$.
  \item Να βρείτε την εξίσωση της εφαπτομένης της $C_f$ η οποία είναι παράλληλη στην ευθεία $y = \dfrac{3}{4}x + 2024$.
  \item Να υπολογίσετε το εμβαδόν του χωρίου που περικλείεται από τη $C_f$, την πλάγια ασύμπτωτή της και τις κατακόρυφες ευθείες $x=1$ και $x=e$.
\end{erwthma}
\end{thema}
%=========== ΤΕΛΟΣ - 63ο ΘΕΜΑ Γ ===========

%=========== 64ο ΘΕΜΑ Γ ===========
\begin{thema}{Γ}
Ένα ταχύπλοο βρίσκεται στο νερό στο σημείο $A$ το οποίο απέχει κατακόρυφη απόσταση $3$ χλμ από την ευθύγραμμη ακτή (έστω σημείο προβολής $O$). Ο οδηγός θέλει να φτάσει σε ένα παραθαλάσσιο χωριό $B$ που βρίσκεται στην ακτή και απέχει $4$ χλμ από το $O$. Η ταχύτητά του στη θάλασσα είναι $15$ χλμ/ώρα και αν τρέξει στην παραλία μπορεί να κινηθεί με $25$ χλμ/ώρα.
\begin{center}
\begin{tikzpicture}[scale=1.5, font=\sffamily]

    % --- Παράμετροι Σχήματος ---
    % yA = 3 km (απόσταση από ακτή)
    % xB = 4 km (απόσταση από Ο)
    % xM = x km (απόσταση Ο-Μ, μεταβλητή)
    \def\yA{3.0}   % y συντεταγμένη του A
    \def\xB{4.0}   % x συντεταγμένη του B
    \def\xM{2.5}   % x συντεταγμένη του M (αντιπροσωπευτική τιμή για το σχήμα)

    % --- Συντεταγμένες ---
    \coordinate (O) at (0, 0);                 % Προβολή του A (0,0)
    \coordinate (A) at (0, \yA);                % Ταχύπλοο (0,3)
    \coordinate (B) at (\xB, 0);                % Χωριό (4,0)
    \coordinate (M) at (\xM, 0);                % Σημείο αποβίβασης (x,0)

    % --- Σχεδίαση Φόντου ---
    % Θάλασσα
    \fill[\xrwma!15] (-1, 0) rectangle (\xB+1, \yA+1);
    % Παραλία / Ακτή
    \fill[secondcolor!15] (-1, -1) rectangle (\xB+1, 0);

    % --- Σχεδίαση Εδάφους & Αντικειμένων ---
    % Ευθύγραμμη Ακτή
    \draw[ultra thick, black] (-1, 0) -- (\xB+1, 0);

    % Διακεκομμένη γραμμή προβολής (AO)
    \draw[dashed, thin, gray] (A) -- (O);
    % Σημείο κάθετης γωνίας στο O
    \draw (0.3, 0) -- (0.3, 0.3) -- (0, 0.3);

    % Πορεία Ταχύπλοου (AM, στο νερό)
    \draw[line width=2pt, \xrwma] (A) -- (M);
    % Πορεία Οδηγού (MB, στην ακτή)
    \draw[line width=2pt, secondcolor] (M) -- (B);

    % --- Σημειώσεις ---
    % Επεξηγηματικό κείμενο για θάλασσα και παραλία
    \node[color=black, font=\bfseries\small] at (2, \yA+0.5) {Θάλασσα};
    \node[color=secondcolor, font=\bfseries\small] at (2, -0.5) {Παραλία};

    % --- Ετικέτες, Διαστάσεις και Ταχύτητα ---
    % Διάσταση 3 km (AO)
    \draw[<->, >=latex] (-0.3, 0) -- (-0.3, \yA) node[midway, left] {3 \si{km}};
    \draw[dashed, gray] (O) -- (-0.8, 0);
    \draw[dashed, gray] (A) -- (-0.8, \yA);

    % Διάσταση 4 km (OB)
    \draw[<->, >=latex] (0, -0.6) -- (\xB, -0.6) node[midway, below] {4 \si{km}};
    \draw[dashed, gray] (O) -- (0, -0.8);
    \draw[dashed, gray] (B) -- (\xB, -0.8);

    % Διάσταση x km (OM)
    \draw[<->, >=latex] (0, -0.3) -- (\xM, -0.3) node[midway, below] {$x$ \si{km}};
    \draw[dashed, gray] (M) -- (\xM, -0.5);

    % Ενδεικτική ετικέτα για την ταχύτητα στο νερό (v_w = 15 km/h)
    \draw[->, >=latex, \xrwma,yshift=2mm] ($([yshift=2mm]A)!0.5!([yshift=2mm]M)$) -- ++(0.4, -0.45) node[right, font=\small] {$\mathbf{v}_w = 15 \si{ km/h}$};

    % Ενδεικτική ετικέτα για την ταχύτητα στην ακτή (v_b = 25 km/h)
    \draw[->, >=latex, secondcolor] ($([yshift=7mm]M)!0.5!(B)$) -- ++(1.0, 0) node[right, font=\small] {$\mathbf{v}_b = 25 \si{ km/h}$};

    % Κουκκίδες Σημείων
    \filldraw (O) circle (1.5pt) node[below left] {$O$};
    \filldraw (A) circle (1.5pt) node[above left] {$A$};
    \filldraw (B) circle (1.5pt) node[above right] {$B$};
    \filldraw (M) circle (1.5pt) node[below] {$M$};

\end{tikzpicture}
\end{center}
\begin{erwthma}
  \item Έστω ότι αποβιβάζεται στο σημείο $M$ της ακτής σε απόσταση $x$ από το $O$. Να εκφράσετε τον συνολικό χρόνο $T(x)$ της διαδρομής $AMB$ συναρτήσει του $x \in [0, 4]$.
  \item Να βρείτε την παράγωγο $T'(x)$ και να λύσετε την εξίσωση $T'(x)=0$.
  \item Ποια είναι η βέλτιστη απόσταση $x$ για να ελαχιστοποιηθεί ο χρόνος ταξιδιού? 
  \item Υπολογίστε το χρόνο σε περίπτωση απευθείας μετάβασης με σκάφος στο $B$ και επιβεβαιώστε ότι υπολείπεται της βέλτιστης.
\end{erwthma}
\end{thema}
%=========== ΤΕΛΟΣ - 64ο ΘΕΜΑ Γ ===========

%=========== 65ο ΘΕΜΑ Γ ===========
\begin{thema}{Γ}
Δίνεται η συνάρτηση $f(x) = e^x + x - 1$.
\begin{erwthma}
  \item Να αποδείξετε ότι η συνάρτηση $f$ αντιστρέφεται και να βρείτε το πεδίο ορισμού της αντίστροφης συνάρτησης $f^{-1}$.
  \item Να υπολογίσετε τις τιμές $f^{-1}(0)$ και $f^{-1}(e)$ και, στη συνέχεια, να λύσετε την ανίσωση $f^{-1}(x) < 1$.
  \item Να αποδείξετε ότι $\dint_0^e f^{-1}(x) \d x = \dint_0^1 x f'(x) \d x$.
  \item Να υπολογίσετε το εμβαδόν του χωρίου που περικλείεται από τη γραφική παράσταση της $f^{-1}$, τον άξονα $x'x$ και την κατακόρυφη ευθεία $x = e$.
\end{erwthma}
\end{thema}
%=========== ΤΕΛΟΣ - 65ο ΘΕΜΑ Γ ===========

%=========== 66ο ΘΕΜΑ Γ ===========
\begin{thema}{Γ}
Θεωρούμε την εξίσωση $e^x = a x$ με παράμετρο $a \in \mathbb{R}$.
\begin{erwthma}
  \item Να δείξετε ότι η συνάρτηση $g(x) = \dfrac{e^x}{x}$ με $x>0$ παρουσιάζει ελάχιστο στο $x=1$ και να βρείτε την ελάχιστη τιμή της.
  \item Αποδείξτε ότι, για $a > e$, η ευθεία $y = a x$ τέμνει την καμπύλη $y = e^x$ σε ακριβώς δύο σημεία στο πρώτο τεταρτημόριο. 
  \item Για $a \le e$, πόσες λύσεις στο $(0, +\infty)$ μπορεί να έχει η εξίσωση $e^x = a x$; Να αιτιολογήσετε την απάντησή σας.
  \item Να υπολογίσετε το όριο $\lim\limits_{x \to 1} \dfrac{g(x) - e}{(x-1)^2}$.
  \item Να αποδείξετε ότι η $g$ είναι κυρτή στο $(0, +\infty)$ και ότι για $1 < \alpha < \beta$ ισχύει $g'(\alpha) < \dfrac{g(\beta) - g(\alpha)}{\beta - \alpha} < g'(\beta)$.
\end{erwthma}
\end{thema}
%=========== ΤΕΛΟΣ - 66ο ΘΕΜΑ Γ ===========

%=========== 67ο ΘΕΜΑ Γ ===========
\begin{thema}{Γ}
Θεωρούμε τη συνάρτηση $f(x) = \dfrac{\ln x}{x^2}$ για $x > 0$.
\begin{erwthma}
  \item Μελετήστε την $f$ ως προς τα ακρότατα και αποδείξτε ότι λαμβάνει μέγιστη τιμή για $x = \sqrt{e}$.
  \item Να βρείτε αν έχει σημείο καμπής.
  \item Να υπολογίσετε το ολοκλήρωμα $E(a)=\dint_1^a f(x) \d x \ \ (a > 1)$ εφαρμόζοντας ολοκλήρωση κατά παράγοντες. 
  \item Υπολογίστε το όριο $\lim\limits_{a \to +\infty}E(a)$.
\end{erwthma}
\end{thema}
%=========== ΤΕΛΟΣ - 67ο ΘΕΜΑ Γ ===========

%=========== 68ο ΘΕΜΑ Γ ===========
\begin{thema}{Γ}
Μια συνάρτηση έχει τον τύπο $f(x) = x^3 - 3a^2 x + 2a^3$, όπου $a > 0$.
\begin{erwthma}
  \item Να βρείτε τα ακρότατα της $f$.
  \item Να βρείτε τις τιμές του $a$ για τις οποίες $\dint_{-1}^{1}{f(x)}=0$.
  \item Να βρείτε το τοπικό μέγιστο της $f$ ως συνάρτηση του $a$ και να υπολογίσετε την ελάχιστη τιμή του.
  \item Υπολογίστε το εμβαδόν μεταξύ της $f(x)$ και του άξονα $x'x$ συναρτήσει του $a$.
\end{erwthma}
\end{thema}
%=========== ΤΕΛΟΣ - 68ο ΘΕΜΑ Γ ===========

%=========== 69ο ΘΕΜΑ Γ ===========
\begin{thema}{Γ}
Έστω η κυβική $f(x) = x^3 - 6x^2 + 9x$.
\begin{erwthma}
  \item Φέρνουμε την εφαπτομένη της $C_f$ στο σημείο με $x_1 = 1$. Αποδείξτε ότι η εφαπτομένη αυτή τέμνει ξανά την καμπύλη της $f$ στο σημείο με $x_2 = 4$.
  \item Υπολογίστε το εμβαδόν που περικλείεται από το γράφημα και την εφαπτομένη.
  \item Να υπολογίσετε το σημείο καμπής $K$ της $f(x)$ και να αποδείξετε ότι η απόσταση $x_K - x_1$ είναι το $\frac{1}{3}$ της συνολικής διαφοράς $x_2 - x_1$.
  \item Υπολογίστε το όριο $\lim\limits_{h \to 0} \dfrac{f(1+3h) - f(1-h)}{h}$.
\end{erwthma}
\end{thema}
%=========== ΤΕΛΟΣ - 69ο ΘΕΜΑ Γ ===========

%=========== 70ο ΘΕΜΑ Γ ===========
\begin{thema}{Γ}
Έστω συνεχής συνάρτηση $f: \mathbb{R} \to \mathbb{R}$ με $\lim\limits_{x \to -\infty} f(x) = -\infty$, $\lim\limits_{x \to +\infty} \dfrac{f(x)}{x} = 2$ και $\lim\limits_{x \to +\infty} (f(x) - 2x) = 3$. Επιπλέον η $C_f$ τέμνει τον άξονα $x'x$ ακριβώς μια φορά στο $x_0 = -1$.
\begin{erwthma}
  \item Να βρείτε, εάν υπάρχει, την ασύμπτωτη της $f(x)$ στο $+\infty$. 
  \item Υπολογίστε το $\lim\limits_{x \to +\infty} \dfrac{f^2(x) - x^2}{x^2 + f(x)}$.
  \item Εξετάστε το πρόσημο της συνάρτησης $f$.
  \item Να βρείτε τα όρια $\lim\limits_{x \to -\infty} \ln(|f(x)|) \cdot e^{f(x)}$ και $\lim\limits_{x \to 0} \dfrac{1}{f(x)}$.
\end{erwthma}
\end{thema}
%=========== ΤΕΛΟΣ - 70ο ΘΕΜΑ Γ ===========

%=========== 71ο ΘΕΜΑ Γ ===========
\begin{thema}{Γ}
Μια συνεχής συνάρτηση $f: \mathbb{R} \to \mathbb{R}$ ικανοποιεί τη σχέση $f^2(x) - 2x f(x) = 1$ για κάθε $x \in \mathbb{R}$, με αρχική συνθήκη $f(0) = 1$.
\begin{erwthma}
  \item Να αποδείξετε ότι ο τύπος της συνάρτησης είναι $f(x) = x + \sqrt{x^2+1}$.
  \item Να αποδείξετε ότι ισχύει η ταυτότητα $f(x) \cdot f(-x) = 1$ για κάθε πραγματικό αριθμό $x$ και να συμπεράνετε ότι $f(x) > 0$.
  \item Να αποδείξετε ότι η συνάρτηση είναι γνησίως αύξουσα στο $\mathbb{R}$.
  \item Να αποδείξτε την ανισότητα $f(x+1) - f(x) < f'(x+1)$ για κάθε $x \in \mathbb{R}$. 
\end{erwthma}
\end{thema}
%=========== ΤΕΛΟΣ - 71ο ΘΕΜΑ Γ ===========

%=========== 72ο ΘΕΜΑ Γ ===========
\begin{thema}{Γ}
Δίνεται η παραγωγίσιμη συνάρτηση $f: \mathbb{R} \to \mathbb{R}$, η οποία ικανοποιεί τη συνθήκη $f'(x) \cdot e^{f(x)} = 2x$ για κάθε $x \in \mathbb{R}$, ενώ $f(0) = 0$.
\begin{erwthma}
  \item Να δείξετε ότι $f(x)=\ln\left(x^2+1\right)$.
  \item Να μελετήσετε την $f(x)$ ως προς την κυρτότητα και τα σημεία καμπής.
  \item Να αποδείξετε ότι ισχύει $\ln(x^2+1) < x$.
  \item Να υπολογίσετε το εμβαδόν του χωρίου που περικλείεται από την εφαπτομένη της $C_f$ στο σημείο $x=0$, την $C_f$ και την ευθεία $x=1$.
\end{erwthma}
\end{thema}
%=========== ΤΕΛΟΣ - 72ο ΘΕΜΑ Γ ===========

%=========== 73ο ΘΕΜΑ Γ ===========
\begin{thema}{Γ}
Δίνεται συνάρτηση $f:(0,+\infty)\to\mathbb{R}$, δύο φορές παραγωγίσιμη που ικανοποιεί την σχέση $x f'(x) - 2f(x) = x^3 e^x$ για κάθε $x > 0$, με $f(1) = e$.
\begin{erwthma}
  \item Αποδείξτε ότι $f(x) = x^2 e^x$.
  \item Εξετάστε τη συνάρτηση ως προς την μονοτονία και τα ακρότατα. Στη συνέχεια εξηγήστε γιατί η συνάρτηση είναι κυρτή στο πεδίο της.
  \item Έστω ένα τυχαίο $x > 0$. Να δείξετε ότι $f(x+2) + f(x) > 2f(x+1)$.
  \item Αποδείξτε την ανισότητα $x^2 e^{x-1} \ge 3x - 2$ για κάθε $x > 0$.
\end{erwthma}
\end{thema}
%=========== ΤΕΛΟΣ - 73ο ΘΕΜΑ Γ ===========

%=========== 74ο ΘΕΜΑ Γ ===========
\begin{thema}{Γ}
Μια παραγωγίσιμη συνάρτηση ικανοποιεί την εξίσωση $x f'(x) + 2 f(x) = x \ln x$ με $x \ge 1$, καθώς και τη σχέση $f(1) = -\dfrac{1}{9}$.
\begin{erwthma}
  \item Να αποδείξετε ότι $f(x) = \dfrac{x}{3}\ln x - \dfrac{x}{9}$.
  \item Να βρείτε την εξίσωση της εφαπτομένης της $C_f$ στο σημείο με τετμημένη $x=e$.
  \item Να υπολογίσετε το όριο $\lim\limits_{x \to +\infty} \dfrac{f(x)}{x}$.
  \item Να αποδείξετε την ανισότητα $\ln(x+1) - \ln x < \dfrac{1}{x}$, για $x \ge 1$.
\end{erwthma}
\end{thema}
%=========== ΤΕΛΟΣ - 74ο ΘΕΜΑ Γ ===========

%=========== 75ο ΘΕΜΑ Γ ===========
\begin{thema}{Γ}
Δίνεται παραγωγίσιμη συνάρτηση για την οποία ισχύει $f(x) \cdot f'(x) = \dfrac{\ln x}{x}$ για κάθε $x>1$, με $f(e) = 1$ και γνωρίζουμε ότι $f(x) > 0$. 
\begin{erwthma}
  \item Να βρείτε, ότι η συνάρτηση είναι η $f(x) = \ln x$ για τα $x \in (1, +\infty)$.
  \item Να μελετήσετε ως προς τη μονοτονία τη συνάρτηση $\varphi(x) = \dfrac{\ln x}{x}$.
  \item Να συγκρίνετε τους αριθμούς $e^\pi$ και $\pi^e$. 
  \item Να αποδείξετε ότι εμβαδόν του χωρίου μεταξύ της $C_f$, του άξονα $x'x$ και της ευθείας $x=e$ ισούται με 1 τ.μ.
\end{erwthma}
\end{thema}
%=========== ΤΕΛΟΣ - 75ο ΘΕΜΑ Γ ===========

%=========== 76ο ΘΕΜΑ Γ ===========
\begin{thema}{Γ}
Δίνεται η παραγωγίσιμη συνάρτηση $f(x)$ με $f(x) > 1$ για κάθε $x \ge 0$. Η $f$ ικανοποιεί τη σχέση $f'(x) = f(x) \ln(f(x))$, με $f(0) = e$.
\begin{erwthma}
  \item Να δείξετε ότι $f(x) = e^{e^x}$.
  \item Να αποδείξετε ότι η $f(x)$ έχει αντίστροφη, έστω $g(x)=f^{-1}(x)$.
  \item Να δείξετε ότι για κάθε $x \ge 0$ ισχύει $e^{e^{x+1}} - e^{e^x} > e^{x+e^x}$.
  \item Να υπολογίσετε το εμβαδόν του χωρίου μεταξύ της $C_{g^{-1}}$, του άξονα $x'x$ και των κατακόρυφων $x=e, x=e^e$.
\end{erwthma}
\end{thema}
%=========== ΤΕΛΟΣ - 76ο ΘΕΜΑ Γ ===========

%=========== 77ο ΘΕΜΑ Γ ===========
\begin{thema}{Γ}
Μια συνεχής στο $\mathbb{R}$ συνάρτηση $f$ ικανοποιεί την εξίσωση $f^2(x) - 4f(x) = x^2 - 4$ για κάθε $x\in\mathbb{R}$, ενώ $f(0)=2$ και $f(x) \ge 2$.
\begin{erwthma}
  \item Να βρείτε, τον τύπο της $f(x)$.
  \item Να εξετάσετε αν η $f(x)$ είναι παραγωγίσιμη στο $x_0=0$.
  \item Να εντοπίσετε γεωμετρικά το σημείο της $C_f$ (με $x > 0$) που ελαχιστοποιεί την απόσταση αυτού από την αρχή των αξόνων $O(0,0)$.
  \item Εφαρμόζεται το Θεώρημα \eng{Rolle} για την $f(x)$ στο διάστημα $[-1, 1]$? Αιτιολογήστε πλήρως με βάση τις προϋποθέσεις του. 
\end{erwthma}
\end{thema}
%=========== ΤΕΛΟΣ - 77ο ΘΕΜΑ Γ ===========

%=========== 78ο ΘΕΜΑ Γ ===========
\begin{thema}{Γ}
Δίνεται παραγωγίσιμη συνάρτηση $f : \mathbb{R} \to \mathbb{R}$ η οποία ικανοποιεί την ισότητα $(x^2+1) f'(x) + 2x f(x) = 2x$ για κάθε $x\in\mathbb{R}$  με αρχική τιμή $f(0) = 0$.
\begin{erwthma}
  \item Δείξτε ότι ο τύπος της συνάρτησης είναι: $f(x) = \dfrac{x^2}{x^2+1}$.
  \item Μελετήστε την $f$ ως προς τη μονοτονία και βρείτε τις ασύμπτωτες.
  \item Δείξτε ότι η ευθεία $y=x$ βρίσκεται πάνω από την $C_f$ για κάθε $x>0$.
  \item Να αποδείξετε ότι η εξίσωση $f(x) = k$ έχει ακριβώς δύο πραγματικές ρίζες για οποιαδήποτε $k \in (0, 1)$.
\end{erwthma}
\end{thema}
%=========== ΤΕΛΟΣ - 78ο ΘΕΜΑ Γ ===========

%=========== 79ο ΘΕΜΑ Γ ===========
\begin{thema}{Γ}
Για μια παραγωγίσιμη συνάρτηση $f$ ισχύει $(x-1) f'(x) - f(x) = (x-1)^2 e^x$  κάθε $x > 1$, ενώ $f(2) = e^2$.
\begin{erwthma}
  \item Αποδείξτε ότι ο τύπος συνάρτησης είναι ο $f(x) = (x-1)e^x$.
  \item Να μελετήσετε την $f$ ως προς τη μονοτονία και να βρείτε το σύνολο τιμών της.
  \item Να αποδείξετε ότι $x e^{x+1} - (x-1)e^x > x e^x$.
  \item Να βρείτε το εμβαδόν μεταξύ της καμπύλης $C_f$, της εφαπτομένης της στο $x=2$, και της κατακόρυφης $x=1$.
\end{erwthma}
\end{thema}
%=========== ΤΕΛΟΣ - 79ο ΘΕΜΑ Γ ===========

%=========== 80ο ΘΕΜΑ Γ ===========
\begin{thema}{Γ}
Μια συνάρτηση $f: \mathbb{R} \to \mathbb{R}$ ικανοποιεί τη σχέση $f'(x) - f(x) = e^x$ και $f(0) = 0$.
\begin{erwthma}
  \item Να αποδείξετε ο τύπος της $f$ είναι $f(x) = x e^x$.
  \item Μελετήστε την $f$ ως προς την κυρτότητα. Σε ποια θέση βρίσκεται το μοναδικό σημείο καμπής?
  \item Να υπολογίσετε το εμβαδόν του χωρίου που περικλείεται από τη γραφική παράσταση της $f$, τον άξονα $x'x$ και την ευθεία $x=1$.
  \item Υπολογίστε το όριο $\lim\limits_{x \to 0} \dfrac{f(x) - \hm x}{x^2}$.
\end{erwthma}
\end{thema}
%=========== ΤΕΛΟΣ - 80ο ΘΕΜΑ Γ ===========

%=========== 81ο ΘΕΜΑ Γ ===========
\begin{thema}{Γ}
Στο παρακάτω σχήμα δίνεται η γραφική παράσταση $C_f$ μιας δύο φορές παραγωγίσιμης συνάρτησης $f$ ορισμένης στο $[-1, 5]$. Στο σημείο $A(2, 3)$ η εφαπτομένη $(\varepsilon)$ έχει κλίση $-1$.

\begin{center}
\begin{tikzpicture}[scale=0.8]
  % grid
  \draw[help lines, gray!30] (-2,-2) grid (6,5);
  % axes
  \draw[-latex] (-2.3,0) -- (6.3,0) node[right] {$x$};
  \draw[-latex] (0,-2.3) -- (0,5.3) node[above] {$y$};
  % ticks
  \foreach \x in {-1,1,2,3,4,5}
    \draw (\x,0.1) -- (\x,-0.1) node[below,font=\footnotesize] {$\x$};
  \foreach \y in {-1,1,2,3,4}
    \draw (0.1,\y) -- (-0.1,\y) node[left,font=\footnotesize] {$\y$};
  % curve (a smooth bell-like shape)
  \draw[very thick, maincolor, smooth, tension=0.6]
    plot coordinates {(-1,-1) (0,0) (1,2.5) (2,3) (3,1) (4,0) (5,-1.5)};
  % tangent line at A(2,3)
  \draw[secondcolor!60!black, thick] (0.5,4.5) -- (5.5,-0.5) node[right,font=\footnotesize] {$(\varepsilon)$};
  % shaded area
  \fill[blue!15, opacity=0.5, smooth, tension=0.6]
    plot coordinates {(0,0) (1,2.5) (2,3) (3,1) (4,0)} -- cycle;
  % points
  \filldraw[maincolor] (2,3) circle (2pt) node[above right] {$A(2,3)$};
  \filldraw[maincolor] (0,0) circle (2pt);
  \filldraw[maincolor] (4,0) circle (2pt);
\end{tikzpicture}
\end{center}

\begin{erwthma}
  \item Να γράψετε την εξίσωση της εφαπτομένης $(\varepsilon)$ της $C_f$ στο σημείο $A(2, 3)$.
  \item Να εξετάσετε αν υπάρχει $\xi \in (0, 4)$ τέτοιο ώστε $f'(\xi) = 0$. Να αιτιολογήσετε.
  \item Να μελετήσετε γραφικά την κυρτότητα της $f$ και να εκτιμήσετε τη θέση ενός σημείου καμπής.
  \item Από το σχήμα, να εκτιμήσετε αν ισχύει $\dint_0^4 f(x) \d x > 4$ ή $\dint_0^4 f(x) \d x < 4$, αιτιολογώντας μέσω σύγκρισης με κατάλληλο ορθογώνιο.
\end{erwthma}
\end{thema}
%=========== ΤΕΛΟΣ - 81ο ΘΕΜΑ Γ ===========

%---------------------------------------------------------------

%=========== 82ο ΘΕΜΑ Γ ===========
\begin{thema}{Γ}
Στο παρακάτω σχήμα δίνεται η γραφική παράσταση της \textbf{παραγώγου} $f'$ μιας συνεχούς και παραγωγίσιμης συνάρτησης $f$ ορισμένης στο $[0, 6]$.

\begin{center}
\begin{tikzpicture}[scale=0.7]
  % grid
  \draw[help lines, gray!30] (-1,-5) grid (7,4);
  % axes
  \draw[-latex] (-1.3,0) -- (7.3,0) node[right] {$x$};
  \draw[-latex] (0,-5.3) -- (0,4.3) node[above] {$y$};
  % ticks
  \foreach \x in {1,2,3,4,5,6}
    \draw (\x,0.1) -- (\x,-0.1) node[below,font=\footnotesize] {$\x$};
  \foreach \y in {-4,-3,-2,-1,1,2,3}
    \draw (0.1,\y) -- (-0.1,\y) node[left,font=\footnotesize] {$\y$};
  % parabola f'(x) = (x-1)(x-5) = x^2-6x+5
  \draw[very thick, secondcolor, domain=0:6, samples=60] plot (\x, {(\x-1)*(\x-5)});
  % shade areas
  \fill[secondcolor!15, opacity=0.5, domain=0:1, samples=30] (0,0) -- plot (\x, {(\x-1)*(\x-5)}) -- (1,0) -- cycle;
  \fill[secondcolor!10, opacity=0.5, domain=1:5, samples=60] (1,0) -- plot (\x, {(\x-1)*(\x-5)}) -- (5,0) -- cycle;
  \fill[secondcolor!15, opacity=0.5, domain=5:6, samples=30] (5,0) -- plot (\x, {(\x-1)*(\x-5)}) -- (6,0) -- cycle;
  % key points
  \filldraw[secondcolor] (1,0) circle (2.5pt) node[above,font=\footnotesize] {$1$};
  \filldraw[secondcolor] (5,0) circle (2.5pt) node[above,font=\footnotesize] {$5$};
  \filldraw[secondcolor] (3,-4) circle (2.5pt) node[below right,font=\footnotesize] {$(3,-4)$};
  % label
  \node[secondcolor] at (5.5,3) {$C_{f'}$};
\end{tikzpicture}
\end{center}

\begin{erwthma}
  \item Να μελετήσετε τη μονοτονία της $f$ στο $[0, 6]$ και να βρείτε τα ακρότατά της.
  \item Να βρείτε τις θέσεις τυχόν σημείων καμπής της $f$.
  \item Χρησιμοποιώντας το σχήμα, να συγκρίνετε τις τιμές $f(5)$ και $f(0)$.
  \item Ποια είναι η κλίση της εφαπτομένης της $C_f$ στο σημείο με τετμημένη $x_0 = 3$; Αν $f(3) = 1$, γράψτε την εξίσωσή της.
\end{erwthma}
\end{thema}
%=========== ΤΕΛΟΣ - 82ο ΘΕΜΑ Γ ===========

%---------------------------------------------------------------

%=========== 83ο ΘΕΜΑ Γ ===========
\begin{thema}{Γ}
Στο παρακάτω σχήμα δίνεται η γραφική παράσταση $C_f$ μιας συνάρτησης $f$ ορισμένης στο $\mathbb{R}^*$. Η $C_f$ παρουσιάζει πλάγια ασύμπτωτη $y = x - 1$ στο $+\infty$. Η $f$ έχει τοπικό ελάχιστο στο σημείο $E(1, 1)$.

\begin{center}
\begin{tikzpicture}[scale=0.7]
  % grid
  \draw[help lines, gray!30] (-3,-3) grid (8,7);
  % axes
  \draw[-latex] (-3.3,0) -- (8.3,0) node[right] {$x$};
  \draw[-latex] (0,-3.3) -- (0,7.3) node[above] {$y$};
  % ticks
  \foreach \x in {-2,-1,1,2,3,4,5,6,7}
    \draw (\x,0.1) -- (\x,-0.1) node[below,font=\footnotesize] {$\x$};
  \foreach \y in {-2,-1,1,2,3,4,5,6}
    \draw (0.1,\y) -- (-0.1,\y) node[left,font=\footnotesize] {$\y$};
  % oblique asymptote
  \draw[secondcolor!80!black, thick] (-2,-3) -- (8,7) node[right,font=\footnotesize] {$y=x-1$};
  \draw[very thick, maincolor, domain=0.18:8, samples=100] plot (\x, {\x - 1 + 1/\x});
  % points
  \filldraw[maincolor] (1,1) circle (2.5pt) node[below right,font=\footnotesize] {$E(1,1)$};
  % Label
  \node[maincolor] at (6,5.8) {$C_f$};
\end{tikzpicture}
\end{center}

\begin{erwthma}
  \item Να υπολογίσετε τα όρια $\displaystyle\lim\limits_{x \to +\infty} f(x)$, $\displaystyle\lim\limits_{x \to +\infty} (f(x) - (x-1))$ και $\displaystyle\lim\limits_{x \to 0^+} f(x)$.
  \item Να μελετήσετε τη μονοτονία της $f$ στο $(0, +\infty)$ και να βρείτε τα ακρότατά της.
  \item Πόσες ρίζες έχει η εξίσωση $f(x) = c$ για τις διάφορες τιμές του $c \in \mathbb{R}$?
  \item Ένα σημείο $M(x, f(x))$ κινείται πάνω στην $C_f$ για $x > 1$ με $x'(t) = 2$ $\mu$/$\sec$. Τη στιγμή $t_0$ που $x = 2$, να βρείτε τον ρυθμό μεταβολής της τεταγμένης.
\end{erwthma}
\end{thema}
%=========== ΤΕΛΟΣ - 83ο ΘΕΜΑ Γ ===========

%---------------------------------------------------------------

%=========== 84ο ΘΕΜΑ Γ ===========
\begin{thema}{Γ}
Στο παρακάτω σχήμα δίνεται η γραφική παράσταση $C_f$ μιας
\textbf{συνεχούς} και \textbf{μη αρνητικής} συνάρτησης $f$,
παραγωγίσιμης στο $(0,4)$, ορισμένης στο $[0,4]$.
Γνωρίζουμε ότι $f'(2)=2$.

\begin{center}
\begin{tikzpicture}[scale=0.9]
  % πλέγμα
  \draw[help lines, gray!30] (-0.5,-0.5) grid (5,5);
  % άξονες
  \draw[-latex] (-0.7,0) -- (5.3,0) node[right] {$x$};
  \draw[-latex] (0,-0.7) -- (0,5.3) node[above] {$y$};
  % υποδιαιρέσεις
  \foreach \x in {1,2,3,4}
    \draw (\x,0.08) -- (\x,-0.08) node[below,font=\footnotesize] {$\x$};
  \foreach \y in {1,2,3,4}
    \draw (0.08,\y) -- (-0.08,\y) node[left,font=\footnotesize] {$\y$};
  % χορδή AΔ: y = x
  \draw[gray!85, thick] (0,0) -- (4,4);
  \node[gray!85, font=\footnotesize] at (1.15,1.5) {$A\varDelta$};
  % S-καμπύλη: f < x στο (0,2), f > x στο (2,4)
  % κλίση ~2 στο K(2,2)
  \draw[very thick, maincolor, smooth, tension=0.6]
    plot coordinates {
      (0,0)
      (0.5,0.04)
      (1,0.25)
      (1.5,0.85)
      (2,2)
      (2.5,3.15)
      (3,3.75)
      (3.5,3.95)
      (4,4)
    };
  % σκίαση κάτω από καμπύλη (ίδια coordinates)
  \fill[maincolor!15, opacity=0.6, smooth, tension=0.6]
    (0,0)
    -- plot coordinates {
         (0,0)(0.5,0.04)(1,0.25)(1.5,0.85)
         (2,2)(2.5,3.15)(3,3.75)(3.5,3.95)(4,4)
       }
    -- (4,0) -- cycle;
  % εφαπτομένη (ε) στο K: y = 2x-2, κλίση 2
  % x=1 -> y=0,  x=3 -> y=4
  \draw[secondcolor!70!black, thick] (1.3,0) -- (2.7,4)
    node[left, font=\footnotesize, secondcolor!70!black] {$(\varepsilon)$};
  % σημεία
  \filldraw[maincolor] (2,2) circle (2.5pt)
    node[left=4pt, font=\small] {$K(2,2)$};
  \filldraw[maincolor] (0,0) circle (2pt)
    node[below left, font=\small] {$A$};
  \filldraw[maincolor] (4,4) circle (2pt)
    node[above right, font=\small] {$\Delta$};
\end{tikzpicture}
\end{center}

\begin{erwthma}

  \item Να μελετήσετε {γραφικά} την κυρτότητα της $f$ στο
    $[0,4]$ και να δικαιολογήσετε ότι το $K(2,2)$ είναι
    {σημείο καμπής} της $C_f$.
  \item Να γράψετε την εξίσωση της εφαπτομένης $(\varepsilon)$ της
    $C_f$ στο $K(2,2)$. Στη συνέχεια να δείξετε ότι υπάρχει $c\in(0,4)$ τέτοιο
    ώστε $f'(c)=1$.
  \item Αιτιολογήστε τις ανισότητες:
      $\dint_0^2 f(x)\,\d x > 0 \text{ και }\dint_2^4 f(x)\,\d x < 8.$
  \item Χρησιμοποιώντας τη θέση της $C_f$ σχετικά με τη
    {χορδή $A\Delta$}, να αποδείξετε ότι:
    \[
      \int_0^2 f(x)\,\d x < 2
      \qquad\text{και}\qquad
      \int_2^4 f(x)\,\d x > 6,
    \]

\end{erwthma}
\end{thema}
%=========== ΤΕΛΟΣ - 84ο ΘΕΜΑ Γ ===========

%---------------------------------------------------------------

%=========== 85ο ΘΕΜΑ Γ ===========
\begin{thema}{Γ}
Στο παρακάτω σχήμα δίνεται η γραφική παράσταση της \textbf{παραγώγου} $f'$ μιας παραγωγίσιμης $f:[-3, 5]\to\mathbb{R}$. Δίνεται ότι $f(-3) = 1$.

\begin{center}
\begin{tikzpicture}[scale=0.65]
  \draw[help lines, gray!30] (-4,-3) grid (6,5);
  % f' segments
  \draw[very thick, secondcolor] (-3,0) -- (-1,4) -- (2,-2) -- (5,1);
  % shade areas
  \fill[secondcolor!15, opacity=0.5] (-3,0) -- (-1,4) -- (-1,0) -- cycle;
  \fill[secondcolor!15, opacity=0.5] (-1,0) -- (-1,4) -- (1,0) -- cycle;
  \fill[secondcolor!10, opacity=0.5] (1,0) -- (2,-2) -- (2,0) -- cycle;
  \fill[secondcolor!10, opacity=0.5] (2,0) -- (2,-2) -- (4,0) -- cycle;
  \fill[secondcolor!15, opacity=0.5] (4,0) -- (5,1) -- (5,0) -- cycle;
  % key points
  \filldraw[secondcolor] (-3,0) circle (2pt);
  \filldraw[secondcolor] (-1,4) circle (2pt) node[above,font=\footnotesize] {$(-1,4)$};
  \filldraw[secondcolor] (2,-2) circle (2pt) node[below,font=\footnotesize] {$(2,-2)$};
  \filldraw[secondcolor] (5,1) circle (2pt) node[above right,font=\footnotesize] {$(5,1)$};
  % zeros of f'
  \filldraw[secondcolor] (1,0) circle (2pt) node[above right,font=\tiny] {$x_1$};
  \filldraw[secondcolor] (4,0) circle (2pt) node[above right,font=\tiny] {$x_2$};
  \node at (5,2) {$C_{f'}$};
\draw[-latex] (-4.3,0) -- (6.3,0) node[right] {$x$};
  \draw[-latex] (0,-3.3) -- (0,5.3) node[above] {$y$};
  \foreach \x in {-3,-2,-1,1,2,3,4,5}
    \draw (\x,0.1) -- (\x,-0.1) node[below,font=\footnotesize] {$\x$};
  \foreach \y in {-2,-1,1,2,3,4}
    \draw (0.1,\y) -- (-0.1,\y) node[left,font=\footnotesize] {$\y$};
\end{tikzpicture}
\end{center}

\begin{erwthma}
  \item Να βρείτε τις ρίζες της $f'$ και να κατασκευάσετε πίνακα μεταβολών της $f$.
  \item Να υπολογίσετε τις τιμές $f(-1)$ και $f(2)$.
  \item Να συγκρίνετε τα $f(-1)$ και $f(5)$, αιτιολογώντας μέσω εμβαδών, και να βρείτε αν η εξίσωση $f(x) = f(-1)$ έχει λύση στο $(2, 5)$.
  \item Σε ποιο σημείο $\xi \in (-3, 5)$ η εφαπτομένη της $C_f$ είναι παράλληλη στη χορδή από $(-3, f(-3))$ στο $(5, f(5))$?
\end{erwthma}
\end{thema}
%=========== ΤΕΛΟΣ - 85ο ΘΕΜΑ Γ ===========

%---------------------------------------------------------------

%=========== 86ο ΘΕΜΑ Γ ===========
\begin{thema}{Γ}
Στο σχήμα δίνεται η γραφική παράσταση $C_f$ μιας δύο φορές παραγωγίσιμης συνάρτησης $f$ ορισμένης στο $[0, 6]$. Είναι γνωστό ότι $f(6) = -2$.

\begin{center}
\begin{tikzpicture}[scale=0.7]
  \draw[help lines, gray!30] (-1,-3) grid (7,4);
  % curve
  \draw[very thick, maincolor, smooth, tension=0.55]
    plot coordinates {(0,0) (0.5,1) (1,2.2) (1.5,2.8) (2,3) (2.5,2.8) (3,2) (3.5,1) (4,0) (5,-1.5) (6,-2)};
  % shaded positive area
  \fill[maincolor!12, opacity=0.5, smooth, tension=0.55]
    (0,0) -- plot coordinates {(0,0) (0.5,1) (1,2.2) (1.5,2.8) (2,3) (2.5,2.8) (3,2) (3.5,1) (4,0)} -- cycle;
  % shaded negative area
  \fill[maincolor!10, opacity=0.5, smooth, tension=0.55]
    (4,0) -- plot coordinates {(4,0) (5,-1.5) (6,-2)} -- (6,0) -- cycle;
  \filldraw[maincolor] (0,0) circle (2pt);
  \filldraw[maincolor] (2,3) circle (2pt) node[above,font=\footnotesize] {$(2,3)$};
  \filldraw[maincolor] (4,0) circle (2pt);
  \filldraw[maincolor] (6,-2) circle (2pt) node[below right,font=\footnotesize] {$(6,-2)$};
  \node at (4.5,2) {$C_f$};
  \draw[-latex] (-1.3,0) -- (7.3,0) node[right] {$x$};
  \draw[-latex] (0,-3.3) -- (0,4.3) node[above] {$y$};
\foreach \x in {1,2,3,4,5,6}
    \draw (\x,0.1) -- (\x,-0.1) node[below,font=\footnotesize] {$\x$};
  \foreach \y in {-2,-1,1,2,3}
    \draw (0.1,\y) -- (-0.1,\y) node[left,font=\footnotesize] {$\y$};
\end{tikzpicture}
\end{center}

\begin{erwthma}
  \item Να μελετήσετε την $f$ ως προς τη μονοτονία, να βρείτε τα τοπικά της ακρότατα και να προσδιορίσετε το σύνολο τιμών της. Στη συνέχεια, να βρείτε το πλήθος των ριζών της εξίσωσης $f(x) = 1$.
  \item Να ελέγξετε αν ικανοποιούνται οι προϋποθέσεις του Θεωρήματος \eng{Rolle} για την $f$ στο διάστημα $[0, 4]$ και να διατυπώσετε το γεωμετρικό συμπέρασμα του θεωρήματος.
  \item Να αποδείξετε γεωμετρικά ότι $\dint_0^4 f(x)\,\d x > 6$.
  \item Αν επιπλέον γνωρίζετε ότι η εφαπτομένη της $C_f$ στο σημείο $O(0,0)$ έχει εξίσωση $y=2x$, να υπολογίσετε το όριο:
  \[ \lim\limits_{x \to 0} \frac{f(x) - \hm x}{x^2 + x} \]
\end{erwthma}
\end{thema}
%=========== ΤΕΛΟΣ - 86ο ΘΕΜΑ Γ ===========

%---------------------------------------------------------------

%=========== 87ο ΘΕΜΑ Γ ===========
\begin{thema}{Γ}
Στο σχήμα δίνονται η γραφική παράσταση $C_f$ μιας γνησίως αύξουσας συνάρτησης $f:[0, 5]\to [0,4]$ και η γραφική παράσταση $C_{f^{-1}}$ της αντίστροφης της.

\begin{center}
\begin{tikzpicture}[scale=0.75]
  \draw[help lines, gray!30] (-1,-1) grid (6,5);
  \draw[-latex] (-1.3,0) -- (6.3,0) node[right] {$x$};
  \draw[-latex] (0,-1.3) -- (0,5.3) node[above] {$y$};
  \foreach \x in {1,2,3,4,5}
    \draw (\x,0.1) -- (\x,-0.1) node[below,font=\footnotesize] {$\x$};
  \foreach \y in {1,2,3,4}
    \draw (0.1,\y) -- (-0.1,\y) node[left,font=\footnotesize] {$\y$};
  \draw[dashed, gray!50] (-1,-1) -- (5,5);
  % f: concave, below y=x
  \node[secondcolor!80!black] at (1,3.5) {$C_{f^{-1}}$};
  \filldraw[maincolor] (0,0) circle (2pt);
  \filldraw[maincolor] (1,1) circle (2pt) node[below right,font=\footnotesize] {$(1,1)$};
  \filldraw[maincolor] (3,2) circle (2pt) node[below right,font=\footnotesize] {$(3,2)$};
  \filldraw[maincolor] (5,4) circle (2pt) node[below right,font=\footnotesize] {$(5,4)$};
  % shaded between f and f^{-1}
  \fill[secondcolor!15, opacity=0.5, smooth, tension=0.4]
    plot coordinates {(1,1) (2,1.6) (3,2) (4,3)}
    -- plot[smooth, tension=0.4] coordinates {(4,3) (3,4) (2,3) (1.6,2) (1,1)} -- cycle;
  \draw[very thick, maincolor, smooth, tension=0.4]
    plot coordinates {(0,0) (1,1) (2,1.6) (3,2) (4,3) (5,4)};
  \node[maincolor] at (4.5,2) {$C_f$};
  % f^{-1}: convex, above y=x
  \draw[very thick, secondcolor!80!black, smooth, tension=0.4]
    plot coordinates {(0,0) (1,1) (1.6,2) (2,3) (3,4) (4,5)};
\end{tikzpicture}
\end{center}

\begin{erwthma}
  \item Να βρείτε τα $f^{-1}(2)$, $f^{-1}(4)$ και $(f \circ f^{-1})(3)$.
  \item Υπολογίστε το $\dint_0^4 f^{-1}(y)\,dy$, αν γνωρίζετε ότι $\dint_0^5 f(x)\,\d x = 9$.
  \item Αν η $f$ είναι κοίλη στο $[0, 5]$, τι κυρτότητα έχει η $f^{-1}$; Αιτιολογήστε γραφικά και αλγεβρικά.
  \item Ένα σημείο κινείται στην $C_{f^{-1}}$ με τη τεταγμένη να αυξάνεται με ρυθμό $3$ $\mu/\sec$. Τη στιγμή που η τεταγμένη είναι $y_0 = 3$, να βρείτε πόσο γρήγορα μεταβάλλεται η τετμημένη.
\end{erwthma}
\end{thema}
%=========== ΤΕΛΟΣ - 87ο ΘΕΜΑ Γ ===========

%---------------------------------------------------------------

%=========== 88ο ΘΕΜΑ Γ ===========
\begin{thema}{Γ}
Στο σχήμα δίνεται η γραφική παράσταση της \textbf{δεύτερης παραγώγου} $f''$ μιας δύο φορές παραγωγίσιμης συνάρτησης $f:[0, 6]\to\mathbb{R}$. Δίνεται $f'(0) = 0$ και $f(0) = 1$.

\begin{center}
\begin{tikzpicture}[scale=0.65]
  \draw[help lines, gray!30] (-1,-4) grid (7,3);
  \draw[-latex] (-1.3,0) -- (7.3,0) node[right] {$x$};
  \draw[-latex] (0,-4.3) -- (0,3.3) node[above] {$y$};
  \foreach \x in {1,2,3,4,5,6}
    \draw (\x,0.1) -- (\x,-0.1) node[below,font=\footnotesize] {$\x$};
  \foreach \y in {-3,-2,-1,1,2}
    \draw (0.1,\y) -- (-0.1,\y) node[left,font=\footnotesize] {$\y$};
  % f'' curve
  \draw[very thick, secondcolor, smooth, tension=0.5]
    plot coordinates {(0,2) (1,1) (2,0) (3,-2) (4,-3) (5,-2) (6,-1)};
  \filldraw[secondcolor] (0,2) circle (2pt) node[above right,font=\footnotesize] {$(0,2)$};
  \filldraw[secondcolor] (2,0) circle (2pt) node[above right,font=\footnotesize] {$(2,0)$};
  \filldraw[secondcolor] (4,-3) circle (2pt) node[below right,font=\footnotesize] {$(4,-3)$};
  \node[secondcolor] at (6,1.5) {$C_{f''}$};
\end{tikzpicture}
\end{center}

\begin{erwthma}
  \item Να μελετήσετε την κυρτότητα της $f$ και να βρείτε σημεία καμπής.
  \item Χρησιμοποιώντας το πρόσημο $f''$ και δεδομένου ότι $f'(0) = 0$, να εξηγήσετε γιατί η $f'$ είναι γνησίως αύξουσα στο $[0, 2]$ και γνησίως φθίνουσα στο $(2, 6)$. Τι συμπεραίνετε για τη μονοτονία της $f$?
  \item Αν $f'$ παίρνει ολικό μέγιστο στο $x = 2$, να αποδείξετε ότι $f'(2) > 0$.
  \item Εξετάστε αν η $f$ παρουσιάζει τοπικό μέγιστο σε κάποιο $\xi > 2$.
\end{erwthma}
\end{thema}
%=========== ΤΕΛΟΣ - 88ο ΘΕΜΑ Γ ===========

%---------------------------------------------------------------

%=========== 89ο ΘΕΜΑ Γ ===========
\begin{thema}{Γ}
Δίνονται οι γραφικές παραστάσεις δύο συνεχών παραγωγίσιμων συναρτήσεων $f$ και $g$ στο $[0, 4]$. Θεωρούμε $h(x) = f(g(x))$.

\begin{center}
\begin{tikzpicture}[scale=0.7]
  \draw[help lines, gray!30] (-1,-1) grid (5,4);
  \draw[-latex] (-1.3,0) -- (5.3,0) node[right] {$x$};
  \draw[-latex] (0,-1.3) -- (0,4.3) node[above] {$y$};
  \foreach \x in {1,2,3,4}
    \draw (\x,0.1) -- (\x,-0.1) node[below,font=\footnotesize] {$\x$};
  \foreach \y in {1,2,3}
    \draw (0.1,\y) -- (-0.1,\y) node[left,font=\footnotesize] {$\y$};
  % f
  \draw[very thick, maincolor, smooth, tension=0.5]
    plot coordinates {(0,1) (1,2.2) (2,3) (3,2.8) (4,2)};
  \node[maincolor] at (0.5,2.8) {$C_f$};
  % g
  \draw[very thick, secondcolor!50!black, smooth, tension=0.5]
    plot coordinates {(0,0) (0.5,1) (1,2) (2,2.5) (3,1) (4,3)};
  \node[secondcolor!50!black] at (3.5,0.5) {$C_g$};
  % key points
  \filldraw[maincolor] (0,1) circle (2pt);
  \filldraw[maincolor] (2,3) circle (2pt) node[above right,font=\footnotesize] {$(2,3)$};
  \filldraw[secondcolor!50!black] (1,2) circle (2pt) node[right,font=\footnotesize] {$(1,2)$};
  \filldraw[secondcolor!50!black] (3,1) circle (2pt) node[below right,font=\footnotesize] {$(3,1)$};
\end{tikzpicture}
\end{center}

\begin{erwthma}
  \item Να υπολογίσετε τις τιμές $h(0)$, $h(1)$ και $h(3)$.
  \item Να υπολογίσετε $h'(1)$, εκτιμώντας γραφικά τις $f'(2)$ και $g'(1)$.
  \item Να εξετάσετε αν η συνάρτηση $h$ παρουσιάζει τοπικό ακρότατο στο $x = 1$.
  \item Να αποδείξετε ότι υπάρχει $\xi \in (0, 4)$ τέτοιο ώστε $h(\xi) = 2$.
\end{erwthma}
\end{thema}
%=========== ΤΕΛΟΣ - 89ο ΘΕΜΑ Γ ===========

%---------------------------------------------------------------

%=========== 90ο ΘΕΜΑ Γ ===========
\begin{thema}{Γ}
Στο σχήμα δίνεται η γραφική παράσταση $C_f$ μιας συνεχούς $f:[-1, 5]\to\mathbb{R}$ αποτελούμενη από ένα ημικυκλικό τόξο ακτίνας $2$ και κέντρου $(1, 0)$, και ένα ευθύγραμμο τμήμα από $(3, 0)$ στο $(5, -3)$.

\begin{center}
\begin{tikzpicture}[scale=0.7]
  \draw[help lines, gray!30] (-2,-4) grid (6,3);
  \draw[-latex] (-2.3,0) -- (6.3,0) node[right] {$x$};
  \draw[-latex] (0,-4.3) -- (0,3.3) node[above] {$y$};
  \foreach \x in {-1,1,2,3,4,5}
    \draw (\x,0.1) -- (\x,-0.1) node[below,font=\footnotesize] {$\x$};
  \foreach \y in {-3,-2,-1,1,2}
    \draw (0.1,\y) -- (-0.1,\y) node[left,font=\footnotesize] {$\y$};
  % semicircle center (1,0), radius 2
  \draw[very thick, maincolor] (-1,0) arc (180:0:2);
  % line from (3,0) to (5,-3)
  \draw[very thick, maincolor] (3,0) -- (5,-3);
  % shaded semicircle area
  \fill[maincolor!12, opacity=0.4] (-1,0) arc (180:0:2) -- cycle;
  % shaded triangle area
  \fill[maincolor!10, opacity=0.4] (3,0) -- (5,-3) -- (5,0) -- cycle;
  \filldraw[maincolor] (-1,0) circle (2pt) node[below left,font=\footnotesize] {$(-1,0)$};
  \filldraw[maincolor] (1,2) circle (2pt) node[above,font=\footnotesize] {$(1,2)$};
  \filldraw[maincolor] (3,0) circle (2pt) node[above right,font=\footnotesize] {$(3,0)$};
  \filldraw[maincolor] (5,-3) circle (2pt) node[right,font=\footnotesize] {$(5,-3)$};
  \node[maincolor] at (4.5,1.5) {$C_f$};
\end{tikzpicture}
\end{center}

\begin{erwthma}
  \item Να βρείτε τον τύπο της συνάρτησης $f$ και στη συνέχεια να τη μελετήσετε ως προς τη μονοτονία και τα ακρότατα.
  \item Να υπολογίσετε (αν υπάρχει) το όριο:
  \[ \lim\limits_{x \to 3} \dfrac{f(x)}{x-3} \]
  \item Να αποδείξετε ότι η εξίσωση $4f(x) - \dint_{-1}^3 f(t)\,\d t = 0$ έχει ακριβώς μία ρίζα στο διάστημα $(1, 3)$.
  \item Ένα σημείο $M(x, y)$ κινείται πάνω στο ημικύκλιο της $C_f$ και η τετμημένη του αυξάνεται με σταθερό ρυθμό $1\,\mu/\si{sec}$. Έστω $A(-1, 0)$ και $B(3, 0)$. Να βρείτε τον ρυθμό μεταβολής του εμβαδού του τριγώνου $MAB$, τη χρονική στιγμή που το $M$ διέρχεται από τον άξονα $y'y$.
\end{erwthma}
\end{thema}
%=========== ΤΕΛΟΣ - 90ο ΘΕΜΑ Γ ===========

%---------------------------------------------------------------

%=========== 91ο ΘΕΜΑ Γ ===========
\begin{thema}{Γ}
Στο παρακάτω σχήμα δίνεται η γραφική παράσταση $C_f$ μιας παραγωγίσιμης συνάρτησης $f:[-2, 4]\to\mathbb{R}$.

\begin{center}
\begin{tikzpicture}[scale=0.7]
  \draw[help lines, gray!30] (-3,-2) grid (5,4);
  \draw[-latex] (-3.3,0) -- (5.3,0) node[right] {$x$};
  \draw[-latex] (0,-2.3) -- (0,4.3) node[above] {$y$};
  \foreach \x in {-2,-1,1,2,3,4}
    \draw (\x,0.1) -- (\x,-0.1) node[below,font=\footnotesize] {$\x$};
  \foreach \y in {-1,1,2,3}
    \draw (0.1,\y) -- (-0.1,\y) node[left,font=\footnotesize] {$\y$};
  \draw[very thick, maincolor, smooth, tension=0.6]
    plot coordinates {(-2,0) (-1,1.6) (0,2.6) (0.5,2.9) (1,3) (1.5,2.4) (2,1) (2.5,-0.4) (3,-1) (3.5,-0.2) (4,1)};
  \filldraw[maincolor] (-2,0) circle (2pt);
  \filldraw[maincolor] (1,3) circle (2pt) node[above,font=\footnotesize] {$(1,3)$};
  \filldraw[maincolor] (3,-1) circle (2pt) node[below,font=\footnotesize] {$(3,-1)$};
  \filldraw[maincolor] (4,1) circle (2pt) node[right,font=\footnotesize] {$(4,1)$};
\end{tikzpicture}
\end{center}

\begin{erwthma}
  \item Να μελετήσετε την $f$ ως προς τη μονοτονία, να βρείτε τις θέσεις των τοπικών ακροτάτων και να αποδείξετε ότι η εξίσωση $f(x) = 0$ έχει ακριβώς τρεις ρίζες στο πεδίο ορισμού της.
  \item Να αποδείξετε ότι υπάρχουν $\xi_1 \in (-2, 1)$ και $\xi_2 \in (1, 3)$ τέτοια, ώστε να ισχύει $2f'(\xi_1) + f'(\xi_2) = 0$.
  \item Να υπολογίσετε το όριο $\lim\limits_{x \to 1} \dfrac{\hm(f(x) - 3)}{x^2 - 1}$.
  \item Να αποδείξετε γεωμετρικά ότι $\dint_{-2}^1 f(x)\,\d x > \dfrac{9}{2}$.
\end{erwthma}
\end{thema}
%=========== ΤΕΛΟΣ - 91ο ΘΕΜΑ Γ ===========

%---------------------------------------------------------------

%=========== 92ο ΘΕΜΑ Γ ===========
\begin{thema}{Γ}
Στο σχήμα δίνεται η γραφική παράσταση $C_{f'}$ της \textbf{παραγώγου} μιας $f:[-1, 5]\to\mathbb{R}$.

\begin{center}
\begin{tikzpicture}
\begin{axis}[
    width=8cm,
    height=8cm,
    axis lines = middle,
    xlabel = {$x$},
    ylabel = {$y$},
    xmin=-1.3, xmax=6.3,
    ymin=-3.3, ymax=3.3,
    xtick={-1,1,2,3,4,5},
    ytick={-2,-1,1,2},
    grid = both,
    grid style={help lines, gray!30},
    samples=100,
    % Ρυθμίσεις για τις ετικέτες των αξόνων στις άκρες των βελών
    every axis x label/.style={at={(ticklabel* cs:1)}, anchor=west},
    every axis y label/.style={at={(ticklabel* cs:1)}, anchor=south},
    tick label style={font=\footnotesize}
]

    % Ορισμός της συνάρτησης f(x) = (x-1)*(x-4)
    \def\f{(x-1)*(x-4)}


    % Η κύρια καμπύλη της συνάρτησης
    \addplot [very thick, secondcolor, domain=-1:5] {\f};

    % Σημεία και Ετικέτες
    \addplot [only marks, mark=*, mark size=1.5pt, secondcolor] 
        coordinates {(1,0) (4,0) (2.5, -2.25)};
    
    \node [below right, font=\footnotesize, secondcolor] at (axis cs:2.5, -2.25) {$(2{.}5,-2{.}25)$};
    
    % Ετικέτα καμπύλης
    \node at (axis cs:5,2) {$C_{f'}$};

\end{axis}
\end{tikzpicture}
\end{center}

\begin{erwthma}
  \item Να μελετήσετε τη μονοτονία και τα ακρότατα της $f$.
  \item Σε ποια $x$ η $C_f$ παρουσιάζει σημεία καμπής?
  \item Αν $f(1) = 4$, να δείξετε ότι $f(4) < 4$.
  \item Βρείτε την εξίσωση της εφαπτομένης της $C_f$ στο σημείο της με τετμημένη $x_0 = 2{,}5$, αν γνωρίζουμε ότι $f(2{,}5) = 1$.
\end{erwthma}
\end{thema}
%=========== ΤΕΛΟΣ - 92ο ΘΕΜΑ Γ ===========

%---------------------------------------------------------------

%=========== 93ο ΘΕΜΑ Γ ===========
\begin{thema}{Γ}
Στο σχήμα δίνεται η $C_f$ μιας δύο φορές παραγωγίσιμης $f:[0, 5]\to\mathbb{R}$. Η εφαπτομένη στο $(2, 4)$ είναι οριζόντια. Η $f$ τέμνει τον $x'x$ μόνο στο $x = 5$.

\begin{center}
\begin{tikzpicture}[scale=0.7]
  \draw[help lines, gray!30] (-1,-1) grid (6,5);
  \draw[-latex] (-1.3,0) -- (6.3,0) node[right] {$x$};
  \draw[-latex] (0,-1.3) -- (0,5.3) node[above] {$y$};
  \foreach \x in {1,2,3,4,5}
    \draw (\x,0.1) -- (\x,-0.1) node[below,font=\footnotesize] {$\x$};
  \foreach \y in {1,2,3,4}
    \draw (0.1,\y) -- (-0.1,\y) node[left,font=\footnotesize] {$\y$};
  \draw[very thick, maincolor, smooth, tension=0.5]
    plot coordinates {(0,1) (0.5,2) (1,3.2) (1.5,3.8) (2,4) (2.5,3.8) (3.5,3) (4.5,1.5) (5,0)};
  \draw[dashed, secondcolor!60!black] (0.5,4) -- (3.5,4);
  \filldraw[maincolor] (0,1) circle (2pt) node[left,font=\footnotesize] {$(0,1)$};
  \filldraw[maincolor] (2,4) circle (2pt) node[above,font=\footnotesize] {$(2,4)$};
  \filldraw[maincolor] (5,0) circle (2pt);
\end{tikzpicture}
\end{center}

\begin{erwthma}
  \item Να προσδιορίσετε τα διαστήματα μονοτονίας της $f$ και το είδος του ακροτάτου στο $x_0=2$. Στη συνέχεια, να υπολογίσετε το όριο $\lim\limits_{x \to 2} \dfrac{1}{f(x) - 4}$.
  \item Να βρείτε την εξίσωση της εφαπτομένης της $C_f$ στο σημείο $A(2, 4)$ και αποδείξετε την ανισότητα $\dint_0^5 f(x)\,\d x < 20$.
  \item Να αποδείξετε ότι η εξίσωση $f(x) = x$ έχει ακριβώς μία ρίζα στο διάστημα $(2, 5)$.
  \item Να υπολογίσετε την τιμή του ολοκληρώματος $I = \dint_0^2 x f''(x)\,\d x$.
\end{erwthma}
\end{thema}
%=========== ΤΕΛΟΣ - 93ο ΘΕΜΑ Γ ===========

%---------------------------------------------------------------

%=========== 94ο ΘΕΜΑ Γ ===========
\begin{thema}{Γ}
Στο σχήμα δίνεται η γραφική παράσταση $C_{f'}$ μιας παραγωγίσιμης συνάρτησης $f:[-3, 3]\to\mathbb{R}$.

\begin{center}
\begin{tikzpicture}[scale=0.7]
  \draw[help lines, gray!30] (-4,-3) grid (4,3);
  \draw[-latex] (-4.3,0) -- (4.3,0) node[right] {$x$};
  \draw[-latex] (0,-3.3) -- (0,3.3) node[above] {$y$};
  \foreach \x in {-3,-2,-1,1,2,3}
    \draw (\x,0.1) -- (\x,-0.1) node[below,font=\footnotesize] {$\x$};
  \foreach \y in {-2,-1,1,2}
    \draw (0.1,\y) -- (-0.1,\y) node[left,font=\footnotesize] {$\y$};
  \draw[very thick, secondcolor, smooth, tension=0.5]
    plot coordinates {(-3,1) (-2.5,0.5) (-2,0) (-1,-1.5) (0,0) (1,2) (2,0) (2.5,-1) (3,-2)};
  \filldraw[secondcolor] (-2,0) circle (2pt) node[above left,font=\footnotesize] {$-2$};
  \filldraw[secondcolor] (0,0) circle (2pt) node[above right,font=\footnotesize] {$0$};
  \filldraw[secondcolor] (2,0) circle (2pt) node[above right,font=\footnotesize] {$2$};
  \node[secondcolor] at (3,1.5) {$C_{f'}$};
\end{tikzpicture}
\end{center}

\begin{erwthma}
  \item Μελετήστε την $f$ ως προς την μονοτονία και την κυρτότητα. Ποια είναι τα τοπικά ακρότατα?
  \item Βρείτε τα σημεία καμπής της $C_f$.
  \item Να συγκρίνετε τις τιμές των τοπικών μεγίστων της $f$, εκτιμώντας γεωμετρικά τα εμβαδά των χωρίων που σχηματίζει η $C_{f'}$ με τον άξονα $x'x$.
  \item Αξιοποιώντας το συμπέρασμα του προηγούμενου ερωτήματος, να αποδείξετε την ανισότητα:
  \[ \dint_{-2}^2 x f''(x)\,\d x < 0 \]
\end{erwthma}
\end{thema}
%=========== ΤΕΛΟΣ - 94ο ΘΕΜΑ Γ ===========

%---------------------------------------------------------------

%=========== 95ο ΘΕΜΑ Γ ===========
\begin{thema}{Γ}
Στο σχήμα δίνεται η γραφική παράσταση $C_f$ μιας συνεχούς συνάρτησης $f:[-3, 3]\to\mathbb{R}$ που είναι {άρτια}.

\begin{center}
\begin{tikzpicture}[scale=0.7]
  \draw[help lines, gray!30] (-4,-2) grid (4,4);
  \draw[-latex] (-4.3,0) -- (4.3,0) node[right] {$x$};
  \draw[-latex] (0,-2.3) -- (0,4.3) node[above] {$y$};
  \foreach \x in {-3,-2,-1,1,2,3}
    \draw (\x,0.1) -- (\x,-0.1) node[below,font=\footnotesize] {$\x$};
  \foreach \y in {-1,1,2,3}
    \draw (0.1,\y) -- (-0.1,\y) node[left,font=\footnotesize] {$\y$};
  \draw[very thick, maincolor, smooth, tension=0.5]
    plot coordinates {(-3,-1) (-2,0) (-1,2) (0,3) (1,2) (2,0) (3,-1)};
  \filldraw[maincolor] (0,3) circle (2pt) node[above right,font=\footnotesize] {$(0,3)$};
  \filldraw[maincolor] (2,0) circle (2pt);
  \filldraw[maincolor] (-2,0) circle (2pt);
  \filldraw[maincolor] (3,-1) circle (2pt) node[right,font=\footnotesize] {$(3,-1)$};
  \draw[dashed, gray!50] (0,-2) -- (0,4);
\end{tikzpicture}
\end{center}

\begin{erwthma}
  \item Να αιτιολογήσετε γιατί η συνάρτηση $f$ δεν αντιστρέφεται. Στη συνέχεια, να λύσετε την ανίσωση $f(f(x)) < 0$ για $x \in [0, 3]$.
  \item Θεωρούμε το ορθογώνιο με κορυφές $A(-x, 0)$, $B(x, 0)$, $\Gamma(x, f(x))$ και $\Delta(-x, f(-x))$, όπου $x \in (0, 2)$. Αν η τετμημένη $x$ αυξάνεται με ρυθμό $1 \mu/\si{sec}$, να υπολογίσετε το ρυθμό μεταβολής του εμβαδού του ορθογωνίου τη χρονική στιγμή που $x = 1$, δεδομένου ότι $f'(1) = -2$.
  \item Αν η $f$ είναι δύο φορές παραγωγίσιμη στο $x_0=0$, να αιτιολογήσετε γιατί $f'(0) = 0$. Στη συνέχεια, αν δίνεται ότι $\lim\limits_{x \to 0} \frac{f(x) - 3}{x^2} = -1$, να υπολογίσετε την τιμή $f''(0)$.
  \item Δίνεται η συνάρτηση $g(x) = x^2 + x \dint_{-2}^2 f(t)\,\d t$. Αν γνωρίζετε ότι $\dint_0^2 f(t)\,\d t = \dfrac{7}{2}$, να βρείτε την εξίσωση της εφαπτομένης της $C_g$ στο σημείο της με τετμημένη $x_0 = 1$.
\end{erwthma}
\end{thema}
%=========== ΤΕΛΟΣ - 95ο ΘΕΜΑ Γ ===========

%---------------------------------------------------------------

%=========== 96ο ΘΕΜΑ Γ ===========
\begin{thema}{Γ}
Στο σχήμα δίνεται η γραφική παράσταση $C_f$ μιας γνησίως αύξουσας συνάρτησης $f:[0, 4]\to[0, 3]$ με $f(0)=0$, $f(4)=3$. Δίνεται επίσης ότι $\dint_0^4 f(x)\,\d x = 4$.

\begin{center}
\begin{tikzpicture}[scale=0.8]
  \draw[help lines, gray!30] (-1,-1) grid (5,4);
  \draw[-latex] (-1.3,0) -- (5.3,0) node[right] {$x$};
  \draw[-latex] (0,-1.3) -- (0,4.3) node[above] {$y$};
  \foreach \x in {1,2,3,4}
    \draw (\x,0.1) -- (\x,-0.1) node[below,font=\footnotesize] {$\x$};
  \foreach \y in {1,2,3}
    \draw (0.1,\y) -- (-0.1,\y) node[left,font=\footnotesize] {$\y$};
  \draw[dashed, gray!50] (-1,-1) -- (4,4);
  \draw[very thick, maincolor, smooth, tension=0.4]
    plot coordinates {(0,0) (1,1.5) (2,2.2) (3,2.7) (4,3)};
  \fill[cyan!10, opacity=0.5, smooth, tension=0.4]
    (0,0) -- plot coordinates {(0,0) (1,1.5) (2,2.2) (3,2.7) (4,3)} -- (4,0) -- cycle;
  \node at (2,0.8) {$E = 4$};
  \draw[dashed] (4,0) -- (4,3) -- (0,3);
  \filldraw[maincolor] (0,0) circle (2pt);
  \filldraw[maincolor] (4,3) circle (2pt) node[above right,font=\footnotesize] {$(4,3)$};
  \node[maincolor] at (3,1) {$C_f$};
\end{tikzpicture}
\end{center}

\begin{erwthma}
  \item Δείξτε ότι η $f$ αντιστρέφεται. Ποιο είναι το πεδίο ορισμού της $f^{-1}$?
  \item Υπολογίστε το ολοκλήρωμα $\dint_0^3 f^{-1}(y)\,dy$.
  \item Να υπολογίσετε το όριο $\lim\limits_{x \to 0^+} \left[ x \hm\left(\frac{1}{f(x)}\right) \right]$.
  \item Να αποδείξετε ότι υπάρχει ακριβώς ένα $\xi \in (0, 4)$ τέτοιο ώστε $4 f(\xi) = \dint_0^4 f(x)\,\d x$.
\end{erwthma}
\end{thema}
%=========== ΤΕΛΟΣ - 96ο ΘΕΜΑ Γ ===========

%---------------------------------------------------------------

%=========== 97ο ΘΕΜΑ Γ ===========
\begin{thema}{Γ}
Στο σχήμα δίνεται η $C_f$ μιας $f:[-1, 5]\to\mathbb{R}$ η οποία είναι ορισμένη παντού εκτός από $x=2$, όπου η $f$ παρουσιάζει κατακόρυφη ασύμπτωτη.

\begin{center}
\begin{tikzpicture}[scale=0.65]
  \draw[help lines, gray!30] (-2,-4) grid (6,4);
  \draw[-latex] (-2.3,0) -- (6.3,0) node[right] {$x$};
  \draw[-latex] (0,-4.3) -- (0,4.3) node[above] {$y$};
  \foreach \x in {-1,1,2,3,4,5}
    \draw (\x,0.1) -- (\x,-0.1) node[below,font=\footnotesize] {$\x$};
  \foreach \y in {-3,-2,-1,1,2,3}
    \draw (0.1,\y) -- (-0.1,\y) node[left,font=\footnotesize] {$\y$};
  \draw (2,-4) -- (2,4);
  \node[font=\footnotesize] at (2.8,3.5) {$x=2$};
  \draw[very thick, maincolor, smooth, tension=0.5]
    plot coordinates {(-1,-1) (1,0) (1.5,1.5) (1.7,3) (1.8,4)};
  \draw[very thick, maincolor, smooth, tension=0.5]
    plot coordinates {(2.2,-4) (2.3,-2.5) (3,-1) (4,0) (5,0.5)};
  \filldraw[maincolor] (0,0) circle (2pt);
  \filldraw[maincolor] (4,0) circle (2pt);
  \node[maincolor] at (5,2) {$C_f$};
\end{tikzpicture}
\end{center}

\begin{erwthma}
  \item Να βρείτε τα όρια $\lim\limits_{x \to 2^-} f(x)$ και $\lim\limits_{x \to 2^+} f(x)$ από τη γραφική παράσταση. Στη συνέχεια, να εξετάσετε αν υπάρχει το όριο:
  \[ \lim\limits_{x \to 2} \frac{e^{f(x)} - 1}{e^{f(x)} + 1} \]
  \item Αν η $f$ είναι παραγωγίσιμη στα διαστήματα $(-1, 2)$ και $(2, 5)$, να αποδείξετε ότι υπάρχουν $\xi_1 \in (0, 1)$ και $\xi_2 \in (3, 4)$ τέτοια ώστε:
  \[ f'(\xi_1) \cdot f'(\xi_2) = \frac{3}{2} \]
  \item Να αποδείξετε ότι η εξίσωση $\displaystyle f(x) = \frac{1}{x-2}$ έχει ακριβώς μία ρίζα στο διάστημα $(2, 5)$.
  \item Αν γνωρίζετε ότι η $f$ είναι κυρτή στο διάστημα $(-1, 2)$, να αποδείξετε ότι $\dint_0^1 f(x)\,\d x < \dfrac{3}{4}$.
\end{erwthma}
\end{thema}
%=========== ΤΕΛΟΣ - 97ο ΘΕΜΑ Γ ===========

%---------------------------------------------------------------

%=========== 98ο ΘΕΜΑ Γ ===========
\begin{thema}{Γ}
Στο σχήμα δίνεται η γραφική παράσταση $C_{f'}$ μιας παραγωγίσιμης συνάρτησης $f:[0, 6]\to\mathbb{R}$. Δίνεται $f(0) = 0$.

\begin{center}
\begin{tikzpicture}[scale=0.7]
  \draw[help lines, gray!30] (-1,-3) grid (7,3);
  \draw[-latex] (-1.3,0) -- (7.3,0) node[right] {$x$};
  \draw[-latex] (0,-3.3) -- (0,3.3) node[above] {$y$};
  \foreach \x in {1,2,3,4,5,6}
    \draw (\x,0.1) -- (\x,-0.1) node[below,font=\footnotesize] {$\x$};
  \foreach \y in {-2,-1,1,2}
    \draw (0.1,\y) -- (-0.1,\y) node[left,font=\footnotesize] {$\y$};
  \draw[very thick, secondcolor] (0,2) -- (2,0) -- (4,-2) -- (6,0);
  \fill[secondcolor!12, opacity=0.5] (0,0) -- (0,2) -- (2,0) -- cycle;
  \fill[secondcolor!10, opacity=0.5] (2,0) -- (4,-2) -- (4,0) -- cycle;
  \fill[secondcolor!10, opacity=0.5] (4,0) -- (4,-2) -- (6,0) -- cycle;
  \filldraw[secondcolor] (0,2) circle (2pt) node[above right,font=\footnotesize] {$(0,2)$};
  \filldraw[secondcolor] (2,0) circle (2pt) node[above right,font=\footnotesize] {$2$};
  \filldraw[secondcolor] (4,-2) circle (2pt) node[below,font=\footnotesize] {$(4,-2)$};
  \filldraw[secondcolor] (6,0) circle (2pt) node[above right,font=\footnotesize] {$6$};
  \node[secondcolor] at (5,2) {$C_{f'}$};
\end{tikzpicture}
\end{center}

\begin{erwthma}
  \item Να μελετήσετε τη συνάρτηση $f$ ως προς τη μονοτονία και την κυρτότητα.
  \item Αξιοποιώντας τα εμβαδά των χωρίων που σχηματίζει η $C_{f'}$ με τον άξονα $x'x$, να υπολογίσετε τις τιμές $f(2)$, $f(4)$ και $f(6)$ και να βρείτε τα ολικά ακρότατα της $f$.
  \item Να αποδείξετε ότι ο τύπος της συνάρτησης είναι:
  \[ f(x) = \begin{cases} 
  -\dfrac{1}{2}x^2 + 2x, & 0 \le x \le 4 \\[2ex]
  \dfrac{1}{2}x^2 - 6x + 16, & 4 < x \le 6
  \end{cases} \]
  \item Να υπολογίσετε το εμβαδόν του χωρίου που περικλείεται από την $C_f$ και τον άξονα $x'x$, καθώς και το όριο:
  \[ \lim\limits_{x \to 2} \dfrac{f(x) - 2}{(x-2) f'(x)} \]
\end{erwthma}
\end{thema}
%=========== ΤΕΛΟΣ - 98ο ΘΕΜΑ Γ ===========

%---------------------------------------------------------------

%=========== 99ο ΘΕΜΑ Γ ===========
\begin{thema}{Γ}
Στο σχήμα δίνεται η $C_f$ μιας συνεχούς και γνησίως αύξουσας συνάρτησης $f:[0, +\infty)\to\mathbb{R}$ με $f(0)=0$, $f(1)=2$, $f(3)=3$, $f(6) = 4$, και οριζόντια ασύμπτωτη την ευθεία $y = 5$ στο $+\infty$.

\begin{center}
\begin{tikzpicture}
\begin{axis}[
    width=11cm, height=7.5cm,
    axis lines=middle,
    xmin=-0.5, xmax=9.5,
    ymin=-0.5, ymax=5.8,
    xtick={1,2,3,4,5,6,7,8,9},
    ytick={1,2,3,4,5},
    xlabel={$x$},
    ylabel={$y$},
    grid=both,
    grid style={line width=.1pt, draw=gray!20},
    major grid style={line width=.2pt, draw=gray!40},
    enlargelimits=false,
    ticklabel style={font=\footnotesize},
    samples=100
]

% Asymptote
\addplot[domain=-0.5:9.5, dashed, secondcolor!50, thick] {5} node[above right, pos=0.8, font=\footnotesize] {$y=5$};

% Curve
\addplot[very thick, maincolor, smooth, tension=0.5] coordinates {
    (0, 0)
    (0.2, 0.8)
    (0.5, 1.4)
    (0.8, 1.8)
    (1, 2)
    (1.5, 2.35)
    (2, 2.63)
    (3, 3)
    (4, 3.35)
    (5, 3.68)
    (6, 4)
    (7, 4.30)
    (8, 4.55)
    (9, 4.75)
    (11, 4.95)
};

% Points
\addplot[mark=*, maincolor, only marks, mark size=1.5pt] coordinates {(0,0) (1,2) (3,3) (6,4)};

% Labels for points
\node[maincolor, above left, font=\footnotesize] at (axis cs:1,2) {$(1,2)$};
\node[maincolor, above left, font=\footnotesize] at (axis cs:3,3) {$(3,3)$};
\node[maincolor, below right, font=\footnotesize] at (axis cs:6,4) {$(6,4)$};

% Label for function
\node[maincolor] at (axis cs:5, 3.2) {$C_f$};

\end{axis}
\end{tikzpicture}
\end{center}

\begin{erwthma}
  \item Να υπολογίσετε το όριο:
  \[ \lim\limits_{x \to +\infty} \frac{x f(x) - x^2 \hm\left(\dfrac{1}{x}\right)}{x + 1} \]
  \item Αν η $f$ είναι παραγωγίσιμη, να αποδείξετε ότι υπάρχει ένα τουλάχιστον $\xi \in (0, 3)$ στο οποίο η εφαπτομένη της $C_f$ σχηματίζει γωνία $45^\circ$ με τον άξονα $x'x$.
  \item Αν γνωρίζετε ότι η $f$ είναι κοίλη στο $[0, 3]$, να αποδείξετε ότι $\dint_0^3 f(x)\,\d x > \dfrac{9}{2}$ και στη συνέχεια ότι:
  \[ \dint_0^3 f^{-1}(x)\,\d x < \dint_0^3 f(x)\,\d x \]
  \item Να αποδείξετε ότι η εξίσωση $f(x) + f(x^3) = 6$ έχει ακριβώς μία ρίζα στο διάστημα $(1, 6)$.
\end{erwthma}
\end{thema}
%=========== ΤΕΛΟΣ - 99ο ΘΕΜΑ Γ ===========

%---------------------------------------------------------------

%=========== 100ο ΘΕΜΑ Γ ===========
\begin{thema}{Γ}
Στο σχήμα δίνονται οι γραφικές παραστάσεις $C_1$ και $C_2$ δύο συναρτήσεων, μιας συνεχούς συνάρτησης $f:[-1, 4]\to\mathbb{R}$ και της παραγώγου της $f'$. Γνωρίζουμε ότι η $C_1$ διέρχεται από τα σημεία $(-1, 0)$, $(1, 3)$ και $(4, -1)$, ενώ η $C_2$ διέρχεται από τα σημεία $(-1, 3)$, $(1, 0)$ και $(4, -2)$.

\begin{center}
\begin{tikzpicture}
\begin{axis}[
    width=11cm, height=8.5cm,
    axis lines=middle,
    xmin=-1.5, xmax=4.5,
    ymin=-2.5, ymax=3.5,
    xtick={-1,1,2,3,4},
    ytick={-2,-1,1,2,3},
    xlabel={$x$},
    ylabel={$y$},
    grid=both,
    grid style={line width=.1pt, draw=gray!20},
    major grid style={line width=.2pt, draw=gray!40},
    enlargelimits=false,
    ticklabel style={font=\footnotesize},
    samples=100
]

% C1 (f) on [-1, 1]
\addplot[very thick, maincolor, domain=-1:1] {-0.75*x^2 + 1.5*x + 2.25};
% C1 (f) on [1, 4]
\addplot[very thick, maincolor, domain=1:4] {(2/27)*(x-1)^3 - (2/3)*(x-1)^2 + 3} node[pos=0.8, above right] {$C_1$};

% C2 (f') on [-1, 1]
\addplot[very thick, secondcolor, domain=-1:1] {-1.5*x + 1.5};
% C2 (f') on [1, 4]
\addplot[very thick, secondcolor, domain=1:4] {(2/9)*(x-1)^2 - (4/3)*(x-1)} node[pos=0.8, below left] {$C_2$};

% Points for C1
\addplot[mark=*, maincolor, only marks, mark size=1.5pt] coordinates {(-1,0) (1,3) (4,-1)};
\node[maincolor, above, font=\footnotesize] at (axis cs:1,3) {$(1,3)$};
\node[maincolor, below left, font=\footnotesize] at (axis cs:4,-1) {$(4,-1)$};

% Points for C2
\addplot[mark=*, secondcolor, only marks, mark size=1.5pt] coordinates {(-1,3) (1,0) (4,-2)};
\node[secondcolor, right, font=\footnotesize] at (axis cs:-1,3) {$(-1,3)$};
\node[secondcolor, below right, font=\footnotesize] at (axis cs:4,-2) {$(4,-2)$};

\end{axis}
\end{tikzpicture}
\end{center}

\begin{erwthma}
  \item Να αιτιολογήσετε ποια από τις καμπύλες $C_1, C_2$ παριστάνει τη συνάρτηση $f$ και ποια την παράγωγό της $f'$. Στη συνέχεια να βρείτε τα διαστήματα μονοτονίας και τα ακρότατα της $f$.
  \item Να αποδείξετε ότι η $f$ είναι κοίλη στο $[-1, 4]$ και ότι η εξίσωση $f(x) = 0$ έχει ακριβώς μία ρίζα στο διάστημα $(1, 4)$.
  \item Να αποδείξετε ότι $\dint_{-1}^1 f(x)\,\d x > 3$ και να υπολογίσετε το εμβαδόν του χωρίου που περικλείεται από την $C_2$, τον άξονα $x'x$ και τις ευθείες $x=-1$ και $x=4$.
  \item Να υπολογίσετε το ολοκλήρωμα $\dint_1^4 x f''(x)\,\d x$.
\end{erwthma}
\end{thema}
%=========== ΤΕΛΟΣ - 100ο ΘΕΜΑ Γ ===========
